<<<<<<< HEAD
\def\thoigian{90}%--Thời gian
\de{ĐỀ THI THỬ LẦN 1 NĂM HỌC 2025-2026}{SỞ GIÁO DỤC HƯNG YÊN}



\begin{center}
	\textbf{PHẦN 1 - Câu trắc nghiệm nhiều phương án lựa chọn.}
\end{center}
\setcounter{ex}{0}

\Opensolutionfile{ans}[ans-ABCD]




\begin{ex}%[Đề thi thử môn Toán Sở GD Hưng Yên - lần 1, năm học 2025-2026]%[GV biên soạn: Nhật Thiện - GV phản biện: Lê Minh]%[2H5N1-2]
	Trong không gian với hệ tọa độ $Oxyz$, cho mặt phẳng $(P)$ có phương trình $x+2y-z-2=0$. Vectơ nào sau đây là một vectơ pháp tuyến của mặt phẳng $(P)$?
	\choice
	{\True $\overrightarrow{c}(1;2;-1)$}
	{$\overrightarrow{b}(1;-2;-1)$}
	{$\overrightarrow{d}(1;-2;1)$}
	{$\overrightarrow{a}(1;2;1)$}
	\loigiai{
		Một vectơ pháp tuyến của mặt phẳng $(P)$ là $\overrightarrow{c}(1;2;-1)$.
	}
\end{ex}
\begin{ex}%[Đề thi thử môn Toán Sở GD Hưng Yên - lần 1, năm học 2025-2026]%[GV biên soạn: Nhật Thiện - GV phản biện: Lê Minh]%[1D6H2-1]
	Cho $\log_ab=5$. Giá trị của biểu thức $\log_a\left(a^2b\right)$ bằng
	\choice
	{$20$}
	{$25$}
	{$10$}
	{\True $7$}
	\loigiai{
		Ta có $\log_a\left(a^2b\right)=\log_a a^2+\log_a b=2+\log_ab=2+5=7$.
	}
\end{ex}
\begin{ex}%[Đề thi thử môn Toán Sở GD Hưng Yên - lần 1, năm học 2025-2026]%[GV biên soạn: Nhật Thiện - GV phản biện: Lê Minh]%[2H2H2-4]
	Trong không gian với hệ tọa độ $Oxyz$, khoảng cách giữa hai điểm $A(-2;-3;9)$ và $B(-2;3;1)$ bằng
	\choice
	{$8$}
	{$6$}
	{$100$}
	{\True $10$}
	\loigiai{
		Khoảng cách $AB=\sqrt{(-2+2)^2+(3+3)^2+(1-9)^2}=10$.
	}
\end{ex}
\begin{ex}%[Đề thi thử môn Toán Sở GD Hưng Yên - lần 1, năm học 2025-2026]%[GV biên soạn: Nhật Thiện - GV phản biện: Lê Minh]%[2H5H1-3]
	Trong không gian với hệ tọa độ $Oxyz$, phương trình mặt phẳng $(P)$ đi qua điểm $I(2;-1;3)$ và có véctơ pháp tuyến $\overrightarrow{n}(2;-2;1)$ là
	\choice
	{$2x-y+3z-9=0$}
	{$2x-2y+z+9=0$}
	{$2x-y+3z+9=0$}
	{\True $2x-2y+z-9=0$}
	\loigiai{
		Phương trình mặt phẳng $(P)$ đi qua điểm $I(2;-1;3)$ và có véctơ pháp tuyến $\overrightarrow{n}(2;-2;1)$ là
		$$2(x-2)-2(y+1)+1(z-3)=0\Leftrightarrow 2x-2y+z-9=0.$$
	}
\end{ex}
\begin{ex}%[Đề thi thử môn Toán Sở GD Hưng Yên - lần 1, năm học 2025-2026]%[GV biên soạn: Nhật Thiện - GV phản biện: Lê Minh]%[2D1N2-2]
	\immini[thm]{Cho hàm số $y=ax^3+bx^2+cx+d$ ($a\ne 0$) có đồ thị như hình vẽ bên. Điểm cực tiểu của đồ thị hàm số đã cho là
		\choice
		{$y=3$}
		{$y=-1$}
		{\True $x=1$}
		{$x=-1$}}{\begin{tikzpicture}[scale=1, font=\footnotesize, line join=round, line cap=round, >=stealth]
			\begin{scope}[scale=.8]
				\draw[->] (-2.5, 0) -- (2.5, 0) node[below] {$x$};
				\draw[->] (0, -2) -- (0, 4.2) node[left] {$y$};
				\node[below left] at (0, 0) {$O$};
				\begin{scope}
					\clip (-2,-2) rectangle (2,4);
					\draw[domain=-2.15:2.15] plot (\x, {(\x)^3 - 3*(\x) + 1});
					\draw[dashed] (-1, 0) node[below] {$-1$} -- (-1, 3) -- (0, 3) node[right] {$3$};
					\draw[dashed] (1, 0) node[above] {$1$} -- (1, -1) -- (0, -1) node[left] {$-1$};
				\end{scope}
				\node[right] at (0.05, 1) {$1$};
			\end{scope}
	\end{tikzpicture}}
	\loigiai{
		Dựa vào đồ thị, điểm cực tiểu của hàm số đã cho là $x=1$.
	}
\end{ex}
\begin{ex}%[Đề thi thử môn Toán Sở GD Hưng Yên - lần 1, năm học 2025-2026]%[GV biên soạn: Nhật Thiện - GV phản biện: Lê Minh]%[1H8H6-2]
	\immini[thm]{Cho hình chóp $S.ABC$ có đáy $ABC$ là tam giác vuông tại $B$, $SA\perp (ABC)$, $SA=AB=a$ (hình vẽ minh họa). Tính số đo của góc nhị diện $[S,BC,A]$.
		\choice
		{$60^\circ$}
		{\True $45^\circ$}
		{$30^\circ$}
		{$90^\circ$}}{
		\begin{tikzpicture}[scale=1, font=\footnotesize, line join=round, line cap=round, >=stealth]
			\path (0,0) coordinate (A)	(1,-1) coordinate (B) (4,0) coordinate (C) (A)+(0,3) coordinate (S);
			\draw (S)--(A)--(B)--(C)--cycle (S)--(B);
			\draw[dashed] (A)--(C);
			\foreach \x/\y in {A/180,B/-90,C/0,S/90} \draw[fill=black] (\x) circle (1pt) node[shift=(\y:0.3cm)] {$\x$};
		\end{tikzpicture}
	}
	\loigiai{
		Ta có $BC\perp AB$ và $BC\perp SA$ suy ra $BC\perp SB$.\\
		Mặt khác $AB\perp BC$.\\
		Do đó, $\widehat{SBA}$ là góc phẳng nhị diện $[S,BC,A]$.\\
		Xét tam giác $SAB$ là tam giác vuông cân tại $A$ nên $\widehat{SBA}=45^\circ$.
		Vậy số đo của góc nhị diện $[S,BC,A]$ bằng $45^\circ$.
	}
\end{ex}
%%Câu 7%%
\begin{ex}%[2D1N4-1]%[Dự án C đề thi thử TN Môn Toán NH25-26 - Dot 4 - Hoàng Tuấn]%[Sở Hưng Yên]%[GVPB - Tên GV]		
	Đường tiệm cận đứng của đồ thị hàm số $y = \dfrac{2x + 3}{x - 2}$ có phương trình là
	\choice
	{$y = 2$}
	{\True $x = 2$}
	{$y = -\dfrac{3}{2}$}
	{$x = -\dfrac{3}{2}$}
	\loigiai{
		Đường tiệm cận đứng của đồ thị hàm số $y = \dfrac{2x + 3}{x - 2}$ có phương trình là $x=2$.
	}
\end{ex}

%%Câu 8%%
\begin{ex}%[2D1N1-1]%[Dự án C đề thi thử TN Môn Toán NH25-26 - Dot 4 - Hoàng Tuấn]%[Sở Hưng Yên]%[GVPB - Tên GV]
	Hàm số $y = x^3 - 3x + 2026$ nghịch biến trên khoảng nào dưới đây?
	\choice
	{$(1; +\infty)$}
	{$(-\infty; 1)$}
	{$(-2; 2)$}
	{\True $(-1; 1)$}
	\loigiai{Tập xác định $\mathscr{D}=\mathbb{R}$.\\
		Ta có $y' = 3x^2 - 3$.\\
		$y'=0\Leftrightarrow \hoac{&x=1\\&x=-1.}$\\
		Bảng biến thiên
		\begin{center}
			\begin{tikzpicture}
				\tkzTabInit[nocadre=false,lgt=1.2,espcl=2.5,deltacl=0.6]
				{$x$/0.6, $y’$/0.6, $y$/2}%
				{$-\infty$,$-1$,$1$,$+\infty$}
				\tkzTabLine{,+,0,-,0,+,}
				\tkzTabVar{-/$-\infty$, +/$2028$,-/$2024$,+/$+\infty$}
			\end{tikzpicture}
		\end{center}
		Từ bảng biến thiên ta thấy  hàm số nghịch biến trên khoảng $(-1; 1)$.
	}
\end{ex}

%%Câu 9%%
\begin{ex}%[2D4H2-1]%[Dự án C đề thi thử TN Môn Toán NH25-26 - Dot 4 - Hoàng Tuấn]%[Sở Hưng Yên]%[GVPB - Tên GV]
	Cho hàm số $y = f(x)$ liên tục trên $\mathbb{R}$, thỏa mãn $\displaystyle\int_{17}^{2026} f(x)\, \mathrm{\,d}x = 8$ và $\displaystyle\int_{89}^{2026} f(x)\, \mathrm{\,d}x = 10$. Giá trị của biểu thức $\displaystyle\int_{17}^{89} f(x)\, \mathrm{\,d}x$ bằng
	\choice
	{$18$}
	{$2026$}
	{$2$}
	{\True $-2$}
	\loigiai{
		Ta có
		\begin{eqnarray*}
			\displaystyle\int_{17}^{89} f(x) \mathrm{\,d}x &=& \displaystyle\int_{17}^{2026} f(x) \mathrm{\,d}x + \displaystyle\int_{2026}^{89} f(x) \mathrm{\,d}x \\
			&=& \displaystyle\int_{17}^{2026} f(x) \mathrm{\,d}x - \displaystyle\int_{89}^{2026} f(x) \mathrm{\,d}x\\
			&=&8-10\\
			&=&-2.
		\end{eqnarray*}
		
	}
\end{ex}

%%Câu 10%%
\begin{ex}%[2D4N1-3]%[Dự án C đề thi thử TN Môn Toán NH25-26 - Dot 4 - Hoàng Tuấn]%[Sở Hưng Yên]%[GVPB - Tên GV]
	Họ nguyên hàm của hàm số $f(x) = \sin x$ là
	\choice
	{$-\sin x + C$}
	{\True $-\cos x + C$}
	{$\sin x + C$}
	{$\cos x + C$}
	\loigiai{
		Họ nguyên hàm của hàm số $f(x)=\sin x$ là $$\displaystyle\int \sin x \,\mathrm{\,d}x = -\cos x + C.$$
	}
\end{ex}


%%Câu 11%%
\begin{ex}%[1D5H1-3]%[Dự án C đề thi thử TN Môn Toán NH25-26 - Dot 4 - Hoàng Tuấn]%[Sở Hưng Yên]%[GVPB - Tên GV]
	Kết quả kiểm tra môn Toán của học sinh lớp $12\text{B}$ được thống kê như sau:
	\begin{center}
		\begin{tabular}{|l|c|c|c|c|c|}
			\hline
			Điểm & $[0; 2)$ & $[2; 4)$ & $[4; 6)$ & $[6; 8)$ & $[8; 10)$ \\
			\hline
			Số học sinh & $1$ & $4$ & $15$ & $16$ & $4$ \\
			\hline
		\end{tabular}
	\end{center}
	Số trung bình cộng của mẫu số liệu ghép nhóm trên bằng
	\choice
	{$6{,}1$}
	{$6$}
	{\True $5{,}9$}
	{$5{,}8$}
	\loigiai{
		Giá trị đại diện của các nhóm lần lượt là $c_1 = 1$, $c_2 = 3$, $c_3 = 5$, $c_4 = 7$, $c_5 = 9$.\\
		Tổng số học sinh $n= 1 + 4 + 15 + 16 + 4 = 40$.\\
		Số trung bình cộng là:
		$\overline{x} = \dfrac{1 \cdot 1 + 3 \cdot 4 + 5 \cdot 15 + 7 \cdot 16 + 9 \cdot 4}{40} = \dfrac{236}{40} = 5{,}9$.
	}
\end{ex}

%%Câu 12%%
\begin{ex}%[1D5H2-3]%[Dự án C đề thi thử TN Môn Toán NH25-26 - Dot 4 - Hoàng Tuấn]%[Sở Hưng Yên]%[GVPB - Tên GV]
	Bảng dưới đây biểu diễn mẫu số liệu ghép nhóm về chiều cao (đơn vị: cm) của $40$ cây cà chua trong vườn ươm:
	\begin{center}
		\begin{tabular}{|c|c|c|c|c|}
			\hline
			Nhóm & $[20; 25)$ & $[25; 30)$ & $[30; 35)$ & $[35; 40)$ \\
			\hline
			Tần số & $5$ & $9$ & $25$ & $1$ \\
			\hline
		\end{tabular}
	\end{center}
	Tứ phân vị thứ ba $Q_3$ (đơn vị: cm) của mẫu số liệu ghép nhóm trên bằng
	\choice
	{$33$}
	{\True $33{,}2$}
	{$33{,}1$}
	{$33{,}4$}
	\loigiai{
		Tổng số mẫu $n = 40$.\\
		Do  $\dfrac{3n}{4} =\dfrac{3\cdot 40}{4}= 30$.\\
		Vị trí của tứ phân vị thứ ba là của $Q_3=\dfrac{x_{30}+x_{31}}{2}$.\\
		Giá trị này rơi vào nhóm $[30;35)$. Đây là nhóm chứa $Q_3$.
		Áp dụng công thức ta có
		\begin{eqnarray*}
			Q_3 &=&u_m + \dfrac{\dfrac{3n}{4} - C}{n_m} \cdot (u_{m+1} - u_m)\\
			&=&30 + \dfrac{30 - 14}{25} \cdot 5\\
			&=&30 + \dfrac{16}{5}\\
			&=&30 + 3{,}2\\
			&=&33{,}2.
		\end{eqnarray*}
	}
\end{ex}



\Closesolutionfile{ans}


\begin{center}
	\textbf{PHẦN 2 - Câu trắc nghiệm đúng sai. Trong mỗi ý a,b,c,d ở mỗi câu, thí sinh chọn đúng hoặc sai}
\end{center}
\setcounter{ex}{0}
\Opensolutionfile{ans}[ans-DS]
\begin{ex}%[Đề thi thử môn Toán SGD Hưng Yên Lần 1 năm học 2025-2026]%[Xuan Vy Pham - Lê Minh]%[2D4V3-3]
	Cho hàm số $f(x)= \sin x +2x$. Hàm số $F(x)$ là một nguyên hàm của hàm số $f(x)$ trên $\mathbb{R}$ và thỏa mãn $F(0)=4$.
	\choiceTF
	{\True Tiếp tuyến của đồ thị hàm số $y=F(x)$ tại điểm có hoành độ $x = \pi$ có hệ số góc $k=2\pi$}
	{\True Họ nguyên hàm của hàm số $f(x)$ là $-\cos x+x^2+C$ (với $C$ là hằng số)}
	{$F(x)= -\cos x+x^2+4$}
	{\True Thể tích khối tròn xoay sinh bởi miền phẳng $D$ giới hạn bởi đồ thị hàm số $y=f(x)$, trục hoành và hai đường thẳng có phương trình $x=0$, $x= \dfrac{\pi}{2}$ bằng $31$ (đvtt) (\textit{không làm tròn kết quả của các phép toán trung gian, chỉ làm tròn kết quả phép toán cuối cùng đến hàng đơn vị})}
	\loigiai{
		\begin{itemchoice}
			\itemch  
			Vì $F(x)$ là một nguyên hàm của hàm số $f(x)$ trên $\mathbb{R}$ nên $F'(x)=f(x)$, $\forall x \in \mathbb{R}$.\\
			Tiếp tuyến của đồ thị hàm số $y=F(x)$ tại điểm có hoành độ $x = \pi$ có hệ số góc là $$k=F'(\pi)=f(\pi)=\sin \pi +2 \pi = 2\pi.$$
			\itemch  Họ nguyên hàm của hàm số $f(x)$ là 
			$$F(x)=\displaystyle \int f(x) \mathrm{\,d}x= \displaystyle \int \left(\sin x +2x\right)\mathrm{\,d}x=-\cos x +x^2+C.$$
			\itemch Vì $F(0)=4$ nên $-\cos 0+0^2+C=4 \Rightarrow C=5$.\\
			Vậy $F(x)=-\cos x +x^2+5$.
			\itemch Thể tích khối tròn xoay sinh bởi miền phẳng $D$ giới hạn bởi đồ thị hàm số $y=f(x)$, trục hoành và hai đường thẳng có phương trình $x=0$, $x= \dfrac{\pi}{2}$ là
			$$V=\pi \displaystyle \int \limits_{0}^{\tfrac{\pi}{2}} \left[f(x)\right]^2 \mathrm{\,d}x = \pi \displaystyle \int \limits_{0}^{\tfrac{\pi}{2}} \left( \sin x+2x\right)^2 \mathrm{\,d}x  \approx 31\; \text{(đvtt)}.$$
		\end{itemchoice}
	}
\end{ex}
\begin{ex}%[Đề thi thử môn Toán SGD Hưng Yên Lần 1 năm học 2025-2026]%[Xuan Vy Pham - Lê Minh]%[2D1V5-8]	
	Một hộ gia đình muốn xây dựng một bể chứa nước có dạng hình hộp chữ nhật không nắp, thể tích $V=12$m$^3$. Đáy bể là hình chữ nhật có chiều dài gấp đôi chiều rộng. Giá thuê nhân công và vật liệu xây đáy là $500$ nghìn đồng/m$^2$, xây thành bể là $300$ nghìn đồng/m$^2$. Gọi $x$ là chiều rộng của đáy bể, $h$ là chiều cao của bể ($x>0$, $h>0$, đơn vị: m).
	\choiceTF
	{\True Thể tích của bể nước được tính bằng công thức $V=2x^2h$ (m$^3$)}
	{\True Chiều cao của bể nước tính theo $x$ là $h=\dfrac{6}{x^2}$ (m)}
	{Tổng chi phí xây dựng bể nước là $T(x)=500x^2+\dfrac{10\,800}{x}$ (nghìn đồng)}
	{Tổng chi phí tối thiểu để xây dựng bể là $9\,324$ (nghìn đồng) (\textit{không làm tròn kết quả của các phép toán trung gian, chỉ làm tròn kết quả phép toán cuối cùng đến hàng đơn vị})}
	\loigiai{\begin{center}
			\begin{tikzpicture}[line join=round, line cap=round,>=stealth,font=\footnotesize,scale=1]
				\def\a{3} 
				\def\b{1.8}
				\def\h{3.5}
				\path 	(0:0) coordinate (A)
				++(0:\a) coordinate (D)
				++(-130:\b) coordinate (C)
				($(A)+(C)-(D)$) coordinate (B)
				($(A)+(90:\h)$) coordinate (A')
				($(B)+(90:\h)$) coordinate (B')
				($(C)+(90:\h)$) coordinate (C')
				($(D)+(90:\h)$) coordinate (D');
				\draw[dashed,thick] 	(B)--(A)--(D)	(A)--(A');
				\draw[thick] (C)--(C') 	(D)--(D') 	(B)--(B') 	(B)--(C)--(D) (A')--(B')--(C')--(D')--cycle;
				\foreach \x/\g in {A/180,B/180,C/0,D/0,A'/180,B'/180,C'/0,D'/0}
				\fill[black] 	(\x) circle (1pt)
				($(\g:4mm)+(\x)$) node {$\x$};	
				\draw ($(A)!1/2!(D)$) node[above]{$2x$};
				\draw ($(C)!1/2!(D)$) node[right]{$x$};
				\draw ($(A)!1/2!(A')$) node[right]{$h$};
			\end{tikzpicture}
		\end{center}
		\begin{itemchoice}
			\itemch Thể tích của bể nước được tính bằng công thức $V=x \cdot (2x) \cdot h =2x^2h$ (m$^3$).
			\itemch Ta có $V=12 \Rightarrow 2x^2 h=12 \Rightarrow h =\dfrac{6}{x^2}$ (m).
			\itemch Chi phí để xây dựng đáy bể là $500 \cdot x \cdot 2x  = 1\, 000x^2$ (nghìn đồng).\\
			Chi phí để xây dựng thành bể là $$300  \left(2x  h+ 4x  h\right)=1\,800 xh= 1\,800 x\cdot \dfrac{6}{x^2} = \dfrac{19 \, 800}{x^2} \; \text{(nghìn đồng)}.$$
			Tổng chi phí xây dựng bể nước là $T(x)=1\,000x^2+\dfrac{10\,800}{x}$ (nghìn đồng).
			\itemch
			Ta có $T'(x)=2 \, 000x-\dfrac{10 \, 800}{x^2}$.\\
			$T'(x) = 0 \Leftrightarrow 2 \, 000x-\dfrac{10 \, 800}{x^2} = 0 \Leftrightarrow 2\,000x^3 - {10 \, 800}=0 \Leftrightarrow x = \dfrac{3}{\sqrt[3]{5}}$. 
			\begin{center}
				\begin{tikzpicture}
					\tkzTabInit[nocadre,lgt=1.2,espcl=2.5,deltacl=0.6]
					{$x$ /1.4, $T’(x)$ /0.6, $T(x)$ /2}
					{$0$,$\dfrac{3}{\sqrt[3]{5}}$,$+\infty$}
					\tkzTabLine{ ,-,z,+,}
					\tkzTabVar{+/$+\infty$,-/$T\left(\dfrac{3}{\sqrt[3]{5}}\right)$,+/$+\infty$}
				\end{tikzpicture}
			\end{center}
			Vì $\min \limits_{(0,+\infty)} T(x)=T\left(\dfrac{3}{\sqrt[3]{5}}\right)$ nên
			tổng chi phí tối thiểu để xây dựng bể là $$T \left(\dfrac{3}{\sqrt[3]{5}}\right) \approx 9\,234 \; \text{(nghìn đồng)}.$$
		\end{itemchoice}
	}
\end{ex}


\begin{ex}%[2H2N2-3]%[2H2H2-2]%[2H2H2-4]%[2H5H1-3]%[Dự án C  đề thi thử TN Môn Toán NH25-26 - Đợt 4- Phan Trung Hiếu]%[TT-SGD-Hưng Yên L1 - NH25-26]%[GVPB: Nguyễn Hoài Nam]
	Trong không gian với hệ tọa độ $Oxyz$, cho hình hộp $ABCD.A'B'C'D'$ có tọa độ các đỉnh $A(1;-1;3)$, $B(0;2;4)$, $D(2;-1;1)$ và $A'(0;1;2)$.
	\choiceTF
	{\True Tọa độ của vectơ $\overrightarrow{AD} = (1; 0; -2)$}
	{\True Tọa độ đỉnh $B'(-1; 4; 3)$}
	{Góc giữa hai vectơ $\overrightarrow{AB}$ và $\overrightarrow{AD}$ là góc nhọn}
	{Phương trình mặt phẳng $(CB'D')$ có dạng $ax + by + cz - 7 = 0$. Khi đó: $a + b - c = 10$}
	\loigiai{
		\begin{center}
			\begin{tikzpicture}[font=\footnotesize, thick, >=latex]
				\path 
				(0,0) coordinate (A)
				(4,0) coordinate (D)
				(-1,-1) coordinate (B)
				($(B)+(D)-(A)$) coordinate (C)
				(1,3) coordinate (A')
				($(B)+(A')-(A)$) coordinate (B')
				($(D)+(A')-(A)$) coordinate (D')
				($(C)+(A')-(A)$) coordinate (C')
				;
				\draw (A')--(B')--(C')--(D')--cycle;
				\draw (B)--(C)--(D) (B)--(B') (C)--(C') (D)--(D');
				\draw[dashed] (B)--(A)--(D) (A)--(A');
				\foreach \x/\g in {A/160,B/180,C/-45,D/0, A'/90,B'/180,C'/90,D'/0}
				\fill[black] 	(\x) circle (1pt)
				($(\g:3mm)+(\x)$) node {$\x$};
			\end{tikzpicture}
		\end{center}
		\begin{itemchoice}
			\itemch $\overrightarrow{AD} = (1;0;-2)$.
			\itemch Ta có $\overrightarrow{AB} = \overrightarrow{A'B'} \Leftrightarrow B' = A' + B - A = \left(-1;4;3\right)$. Vậy $B'\left(-1;4;3\right)$.
			\itemch $\overrightarrow{AB} = (-1;3;1)$. Ta có
			\begin{equation*}
				\cos\left(\overrightarrow{AB},\overrightarrow{AD}\right) = \dfrac{\overrightarrow{AB}\cdot\overrightarrow{AD}}{\left|\overrightarrow{AB}\right|.\left|\overrightarrow{AD}\right|}=\dfrac{-1\cdot1+3\cdot0+1\cdot(-2)}{\sqrt{1^2+0^2+(-2)^2}\cdot\sqrt{(-1)^2+3^2+1^2}} = \dfrac{-3}{\sqrt{55}}<0.
			\end{equation*}
			Do đó góc giữa hai vectơ $\overrightarrow{AB}$ và $\overrightarrow{AD}$ là góc tù.
			\itemch $\overrightarrow{BC} = \overrightarrow{AD}\Leftrightarrow C = B + D- A = (1;2;2)$. Vậy $C(1;2;2)$.\\
			Tương tự, $\overrightarrow{AD} = \overrightarrow{A'D'}\Leftrightarrow D'(1;1;0)$.\\
			Ta có $\overrightarrow{CB'}=(-2;2;1)$, $\overrightarrow{CD'}=(0;-1;-2)$.\\
			Véc-tơ pháp tuyến của mặt phẳng $(CB'D')$ là $\overrightarrow{n}=\left[\overrightarrow{CB'},\overrightarrow{CD'}\right] = (-3;-4;2)$.\\
			Khi đó, phương trình mặt phẳng $(CB'D')$ có dạng $-3x-4y+2z+D=0$.\\
			Vì $C\in(C'B'D')$ nên $-3-8+4+D=0\Leftrightarrow D=7$.\\
			Vậy $(CB'D')\colon 3x+4y-2z-7=0$. Suy ra $a=3$, $b=4$, $c=-2$.\\
			Khi đó, $a+b-c=9$.
		\end{itemchoice}
	}
\end{ex}

\begin{ex}%[0D2V2-3]%[Dự án C  đề thi thử TN Môn Toán NH25-26 - Đợt 4- Phan Trung Hiếu]%[TT-SGD-Hưng Yên L1 - NH25-26]%[GVPB: Nguyễn Hoài Nam]
	Một xưởng sơn dự định sản xuất hai loại sơn là sơn nội thất và sơn ngoài trời. Nguyên liệu để sản xuất gồm hai loại A và B với trữ lượng lần lượt là $6$ tấn và $8$ tấn. Để sản xuất một tấn sơn nội thất cần $2$ tấn nguyên liệu A và $1$ tấn nguyên liệu B. Để sản xuất một tấn sơn ngoài trời cần $1$ tấn nguyên liệu A và $2$ tấn nguyên liệu B. Kết quả nghiên cứu thị trường cho thấy nhu cầu sơn nội thất không quá $2$ tấn và không nhiều hơn nhu cầu sơn ngoài trời $1$ tấn. Giá bán một tấn sơn nội thất là $60$ triệu đồng, một tấn sơn ngoài trời là $30$ triệu đồng (giả sử rằng số sơn sản xuất ra đều được bán hết).
	\choiceTF
	{Gọi $x$; $y$ lần lượt là số tấn sơn nội thất và sơn ngoài trời cần sản xuất $x \ge 0$, $y \ge 0$}
	{Biểu thức doanh thu $F(x, y) = 30x + 60y$ (triệu đồng)}
	{\True Doanh thu lớn nhất của xưởng là $180$ triệu đồng}
	{Trong các phương án sản xuất đem lại doanh thu lớn nhất, biết rằng tổng số lượng sơn cả hai loại dự định sản xuất không quá $4{,}5$ tấn. Khi đó lượng sơn nội thất cần sản xuất ít nhất là $1{,}4$ tấn}
	
	\loigiai{
		\begin{itemchoice}
			\itemch Gọi $x$; $y$ lần lượt là số tấn sơn nội thất và sơn ngoài trời cần sản xuất.\\
			Điều kiện: $0\leq x \leq 2$, $y\geq 0$ và $x\leq y+1$.
			\immini[thm]{
				Ta có hệ phương trình $\heva{&0\leq x\leq 2\\&y \leq 0\\&x-y\leq1\\&2x+y\leq 6&\\&x+2y\leq 8\\&x+y\leq4{,}5.}$\\
				Miền nghiệm là miền đa giác $OAFBCDE$.
				\begin{itemize}
					\item Tại $A(0,4)$ ta có $F(0,4) = 120$.
					\item Tại $F\left(\dfrac{3}{2};3\right)$ ta có $F\left(\dfrac{3}{2};3\right) = 180$.
					\item Tại $C(2,2)$, ta có $F(2,2) = 180$.
					\item Tại $D(2,1)$ ta có $F(2,1) = 150$.
					\item Tại $E(1,0)$ ta có $F(1,0) = 60$.
				\end{itemize}
			}{
				\begin{tikzpicture}[font=\footnotesize, line cap=round, line join = round, thick, >=stealth, scale=.7,transform shape]
					\def\xmin{-1}
					\def\xmax{9}
					\def\ymin{-1}
					\def\ymax{9}
					\clip (\xmin,\ymin-.5) rectangle (\xmax,\ymax);
					\draw[->] (\xmin,0)--(0,0) node[below left]{$O$} -- (\xmax,0)node[above left]{$x$};
					\draw[->] (0,\ymin-.5)--(0,\ymax)node[below right]{$y$};
					\draw (2,\ymin-.5)--(2,\ymax);
					\draw[domain = \xmin:\xmax] plot (\x, {\x -1});
					\draw[domain = \xmin:\xmax] plot (\x, {-2*\x +6});
					\draw[domain = \xmin:\xmax] plot (\x, {-\x/2 +4});
					\draw[domain = \xmin:\xmax] plot (\x, {-\x +4.5});
					\fill[pattern= north west lines, pattern color =  cyan] (\xmin,0) rectangle (\xmax,\ymin-.5);
					\fill[pattern= north west lines, pattern color =  red] (\xmin,\ymin-.5) rectangle (0,\ymax);
					\fill[pattern= north east lines, pattern color =  blue] (2,\ymin-.5) rectangle (\xmax,\ymax);
					\fill[pattern= north east lines, pattern color =  yellow] (2,\ymin-.5) rectangle (\xmax,\ymax);
					\fill[pattern= north west lines, pattern color =  pink] (\xmin,\ymax) -- plot[domain = \xmin:\xmax] (\x, {-\x/2 +4}) -- (\xmax,\ymax);
					%				\fill[pattern= north east lines, pattern color =  magenta!30] (\xmin,\ymax) -- plot [domain = \xmin:\xmax](\x, {-2*\x +6}) -- (\xmax,\ymax);
					\fill[pattern= north east lines, pattern color =  magenta!30] (\xmin,\ymax) -- plot [domain = \xmin:\xmax](\x, {-\x +4.5}) -- (\xmax,\ymax);
					\foreach \x/\y/\pos/\name in {0/4/200/A,1.5/3/-120/B,2/2/0/C,2/1/-45/D,1/0/-90/E, 1/3.5/-120/F}
					\fill[black] 	(\x,\y) circle (1.5pt)
					($(\pos:3mm)+(\x,\y)$) node {$\name$};
				\end{tikzpicture}
			}
			\itemch Biểu thức doanh thu $F(x;y)=60x+30y$ (triệu đồng).
			%			\itemch Các phương án sản xuất đem lại doanh thu lớn nhất biết rằng tổng số lượng sơn cả hai loại dự định sản xuất không quá $4{,}5$ tấn khi sản xuất $1{,}5$ tấn sơn nội thất và $3$ tấn sơn ngoài trời. Khi đó, lượng sơn nội thất cần sản xuất ít nhất là $1{,}5$ tấn.
			\itemch Doanh thu lớn nhất của xưởng là $180$ triệu đồng.
			\itemch Doanh thu đạt giá trị lớn nhất tại điểm $F\left(\dfrac{3}{2};3\right)$ hoặc $C(2,2)$, khi đó lượng sơn nội thất cần sản xuất ít nhất là $1{,}5$ tấn. 
		\end{itemchoice}
		
	}
\end{ex}


\Closesolutionfile{ans}


\begin{center}
	\textbf{PHẦN 3 - Câu trắc nghiệm trả lời ngắn}
\end{center}
\setcounter{ex}{0}

\Opensolutionfile{ans}[ans-KQ]
\begin{ex}%[2D1H2-7]%[Dự án đề TT Toán NH 2025-2026 Đợt 4-Lê Minh Thiện Anh]%[Sở GD Hưng Yên -Lần 1 - Hưng Yên]%[GVPB: Lê Minh]
	Trong trò chơi mô phỏng xây dựng công viên giải trí Roller Coaster Tycoon, một kỹ sư thiết kế một đoạn đường ray hình lượn sóng. Trong hệ trục tọa độ $Oxy$, đoạn đường ray này được mô hình hóa bởi đồ thị của hàm số: $y = x^3 - 6x^2 + 9x + 1$, với $0 \le x \le 4$. Biết hệ trục tọa độ được thiết lập sao cho trục $Ox$ nằm ngang và mỗi đơn vị độ dài trên các trục tọa độ tương ứng với $2$ mét trên thực tế. Để kiểm tra độ ổn định của cấu trúc, kỹ sư sử dụng thiết bị laser để đo khoảng cách trực tiếp giữa hai điểm chốt kỹ thuật đặt tại điểm cực đại $A$ và điểm cực tiểu $B$ của đồ thị hàm số trên. Hãy tính độ dài thực tế của khoảng cách giữa hai điểm $A$ và $B$ \textit{(làm tròn kết quả cuối cùng đến hàng phần trăm theo đơn vị mét)}.
	\shortans[oly]{8{,}94}
	\loigiai{
		Ta có $y' = 3x^2 - 12x + 9$.\\
		$y' = 0 \Leftrightarrow 3(x^2 - 4x + 3) = 0 \Leftrightarrow \hoac{ &x = 1 \Rightarrow y = 5 \\ &x = 3 \Rightarrow y = 1.}$\\
		Tọa độ hai điểm cực trị là $A(1; 5)$ và $B(3; 1)$.\\
		Khoảng cách giữa hai điểm $A$, $B$ trên mặt phẳng tọa độ là
		$$AB = \sqrt{(3-1)^2 + (1-5)^2} = \sqrt{2^2 + (-4)^2} = \sqrt{20} = 2\sqrt{5}.$$
		Vì mỗi đơn vị độ dài tương ứng với $2$ mét thực tế, nên độ dài thực tế là
		$$L = 2\sqrt{5} \cdot 2 = 4\sqrt{5} \approx 8{,}94 \,(\text{m}).$$
	}
\end{ex}

\begin{ex}%[1D6H3-5]%[Dự án đề TT Toán NH 2025-2026 Đợt 4-Lê Minh Thiện Anh]%[Sở GD Hưng Yên -Lần 1 - Hưng Yên]%[GVPB: Lê Minh]
	Chỉ số giá tiêu dùng CPI là chỉ số tính theo phần trăm phản ánh mức thay đổi của giá hàng hóa tiêu dùng theo thời gian. Giả sử tại một quốc gia, chỉ số CPI sau $n$ năm được dự báo theo công thức: $CPI_n = CPI_0(1+g)^n$, trong đó $CPI_0$ là chỉ số tại thời điểm bắt đầu dự báo và $g$ là tỉ lệ lạm phát trung bình hằng năm. Biết rằng vào tháng 1 năm 2020 chỉ số CPI của quốc gia này là $118{,}09$ và tỉ lệ lạm phát duy trì ổn định ở mức $g = 3{,}21\%$/năm. Hãy tính chỉ số CPI dự báo của quốc gia này vào tháng 1 năm $2030$ \textit{(làm tròn kết quả cuối cùng đến hàng đơn vị)}.
	\shortans[oly]{162}
	\loigiai{
		Thời gian từ tháng $1$ năm $2020$ đến tháng $1$ năm $2030$ là $n = 10$ năm.\\
		Áp dụng công thức với $CPI_0 = 118{,}09$, $g = 3,21\% = 0{,}0321$ và $n = 10$, ta có\\
		$$CPI_{10} = 118{,}09 \cdot \left(1 + 0{,}0321\right)^{10} \approx 162.$$
	}
\end{ex}

\begin{ex}%[2D4C3-2]%[TT lần 1, Sở GD-ĐT Hưng Yên, 2526]%[Huỳnh Xuân Tín]
	Để chuẩn bị cho lễ hội, một đơn vị thi công dựng một cổng chào dạng vòm Parabol có chiều cao $4$ m và chiều rộng chân cổng là $4$ m. Ở chính giữa cổng, người ta thiết kế một lối đi hình chữ nhật cao $2$ m và rộng $2$ m. Phần diện tích mặt trước của cổng Parabol (sau khi đã trừ lối đi hình chữ nhật) được chia thành $8$ lớp ngang, mỗi lớp cao $0{,}5$ m để ốp tấm nhôm Alu trang trí có màu khác nhau (như bản vẽ thiết kế hình bēn dưới).
	\begin{center}
		\begin{tikzpicture}[line cap=butt,line join=miter,>=stealth]
			\tikzset{declare function={xmin=-1.5;xmax=3.5;ymin=-4.5;ymax=2;
					f(\x)=-0.5*(\x)^2+8;
				},
				smooth,samples=450
			}
			\draw[->] (-5,0)--(5,0) node[shift={(-100:7pt)},font=\normalsize]{$ x $};
			%	\draw[->] (0,ymin)--(0,ymax) node[shift={(190:7pt)},font=\normalsize]{$ y $};
			%	\fill (0,0) node[shift={(225:7pt)},font=\normalsize]{$ O $};
			\fill (-1,0) circle (1pt); \fill (1,-2) circle (1pt); \fill (2,0) circle (1pt); \fill (3,-4) circle (1pt); \fill (0,-4) circle (1pt);
			\draw (0,7)node[above]{Lớp $8$} (0,6)node[above]{Lớp $7$} 
			(0,5)node[above]{Lớp $6$} 
			(0,4)node[above]{Lớp $5$} 
			(-2,0.5)node[left]{Lớp $1$}  	
			(-2,1.5)node[left]{Lớp $2$} 
			(-2,2.5)node[left]{Lớp $3$} 
			(-2,3.5)node[left]{Lớp $4$} 
			(0,-0.3)node[below]{$2$ m} 
			(0,-0.8)node[below]{$4$ m}
			(0,-1.4)node[below]{Chiều rộng} 
			(-5.5,4)node[left]{$4$m} 
			(-4.5,0.5)node[left]{$0{,}5$m} 
			;
			%	\clip (xmin,ymin) rectangle (xmax,ymax);
			\draw  plot[domain=-4:4] (\x, {f(\x)});
			\draw  (-2,0)--(-2,4)--(2,4)--(2,0) (-3.741657,1)--(-2,1) (3.741657,1)--(2,1) (-3.464101615,2)--(-2,2) (3.464101615,2)--(2,2) (-3.16227766,3)--(-2,3) (3.16227766,3)--(2,3) (-2.828427125,4)--(2.828427125,4) (-2.449489743,5)--(2.449489743,5) (-2,6)--(2,6) (-1.41421562,7)--(1.414213562,7);	
			\draw [<->] (-2,-0.3)--(2,-0.3);
			\draw [<->] (-4,-0.8)--(4,-0.8);
			\draw [<->] (-5.7,0)--(-5.7,8);
			\draw [<->] (-4.3,0)--(-4.3,1);
		\end{tikzpicture}	
	\end{center}
	Đơn giá thi công lắp tấm nhôm Alu cho cồng được tính như sau
	\begin{itemize}
		\item  Lớp dưới cùng có giá là $400$ nghìn đồng/m$^2$.
		\item  Càng lên cao, giá thi công mỗi lớp tiếp theo tăng thêm $50$ nghìn đồng/m$^2$ so với lớp ngay dưới nó. 
	\end{itemize}
	Tính tổng chi phí (đơn vị: nghìn đồng) ốp Alu cho phần diện tích còn lại của cổng sau khi đã loại bỏ lối đi (\textit{không làm tròn kết quá các phép toán trung gian, chỉ làm tròn kết quá phép toán cuối cùng đến hàng đơn vị})?
	\shortans[oly]{3814}
	\loigiai{Gắn hệ trục tọa độ $Oxy$ sao cho đỉnh Parabol nằm trên trục $Oy$, trục $Ox$ đi qua chân vòm. Phương trình Parabol là $(P)\colon y=-x^2+4$.
		\begin{center}
			\begin{tikzpicture}[line cap=butt,line join=miter,>=stealth]
				\tikzset{declare function={xmin=-1.5;xmax=3.5;ymin=-4.5;ymax=2;
						f(\x)=-0.5*(\x)^2+8;
					},
					smooth,samples=450
				}
				\draw[->] (-5,0)--(5,0) node[shift={(-100:7pt)},font=\normalsize]{$ x $};
				\draw[->] (0,-1)--(0,9) node[shift={(190:7pt)},font=\normalsize]{$ y $};
				\fill (0,0) node[shift={(225:7pt)},font=\normalsize]{$ O $};
				\fill (-4,0) circle (1pt)node[below]{$-4$};
				\fill (4,0) circle (1pt)node[below]{$4$};
				\fill (-2,0) circle (1pt)node[below]{$-2$};
				\fill (2,0) circle (1pt)node[below]{$2$};			
				\fill (0,8) circle (1pt)node[above right]{$4$};		
				\clip (-6,-1) rectangle (6,9);
				\draw  plot[domain=-4:4] (\x, {f(\x)});
				\draw  (-2,0)--(-2,4)--(2,4)--(2,0) (-3.741657,1)--(-2,1) (3.741657,1)--(2,1) (-3.464101615,2)--(-2,2) (3.464101615,2)--(2,2) (-3.16227766,3)--(-2,3) (3.16227766,3)--(2,3) (-2.828427125,4)--(2.828427125,4) (-2.449489743,5)--(2.449489743,5) (-2,6)--(2,6) (-1.41421562,7)--(1.414213562,7);	
			\end{tikzpicture}	
		\end{center}
		Diện tích chỏm Parabol được tính nhanh bằng công thức $S=\dfrac{2}{3} B h$ (trong dó $B$ là độ rộng đáy, $h$ là chiều cao).\\ 
		Gọi $S_0, S_1, \ldots, S_8$ là diện tích hình phẳng giới hạn bởi $(P)$ và các đường thẳng $y=y_i$. \\
		Chiều cao từ đỉnh Parabol xuống đường $y_i$ là $h=4-y_i$. \\
		Từ phương trình $(P) \Rightarrow x=\sqrt{4-y} \Rightarrow B=2 x=2 \sqrt{4-y_i}$. \\
		Suy ra diện tích chỏm phía trên đường $y=y_i$ là	
		$$S_i=\dfrac{2}{3} \cdot 2 \sqrt{4-y_i} \cdot\left(4-y_i\right)=\dfrac{4}{3}\left(\sqrt{4-y_i}\right)^3.
		$$
		Ta có $S_0(y=0)$; $S_1(y=0,5)$; $S_2(y=1)$; $\ldots$; $S_8(y=4)$. \\
		Diện tích các lớp ngang (chưa trừ phần lõi) là $L_1=S_0-S_1$; $L_2=S_1-S_2$; $\ldots$; $L_8=S_7-S_8$.\\
		Tổng chi phí cho toàn bộ diện tích Parabol (giả sử chưa trừ lõi) là
		\begin{eqnarray*}
			T_1	&=& 400 L_1+450 L_2+\ldots+750 L_8 \\
			&=&400\cdot\left(S_0-S_1\right)+450\cdot\left(S_1-S_2\right)+\ldots+750\cdot\left(S_7-S_8\right)\\
			&=&400\cdot S_0+50\cdot\left(S_1+S_2+\ldots+S_7\right) \\
			&=&400 \cdot \dfrac{4}{3} \cdot\left(\sqrt{4}\right)^3+50 \displaystyle\sum\limits_{y=0{,}5 \text { step } 0{,}5}^{3{,}5} \dfrac{4}{3}\left(\sqrt{4-y}\right)^3.
		\end{eqnarray*}
		Chi phí cho phần lõi (hình chữ nhât $2$ m $\times 2$ m, nằm gọn trong các lớp $L_1$ đến $L_4$, mỗi lớp phần lõi có diện tích $2 \times 0{,}5=1$ m$^2$)
		$$T_2=400 \cdot 1+450 \cdot 1+500 \cdot 1+550 \cdot 1=1\,900.$$
		Vậy tổng chi phí thực tế là
		$$T=T_1-T_2 \approx 3\,814\, \text {(nghìn đồng)}.$$	
	}
\end{ex}


\begin{ex}%[2D1V4-3]%[TT lần 1, Sở GD-ĐT Hưng Yên, 2526]%[Huỳnh Xuân Tín]%[GVPB: Lê Minh]
	Một quần thể vi khuẩn được nuôi cấy trong phòng thí nghiệm. Nồng độ dinh dưỡng $S$ (đơn vị mg/l) thay đổi theo thời gian $t$ giờ $(t \geq 0)$ được mô hình hóa bởi hàm số $S(t)=\dfrac{10 t+5}{t+1}$. Biết rằng tốc độ sinh trưởng $V$ của vi khuẩn phụ thuộc vào nồng độ dinh dưỡng theo hàm số $V(S)=\dfrac{5 S}{S+2}$. Khi thời gian $t$ kéo dài, tốc độ sinh trưởng $V$ tăng dần và ổn định quanh một ngưỡng $K$ nhất định. Hỏi sau bao nhiêu \textbf{phút} thì tốc độ sinh trưởng của vi khuẩn đạt $90\%$ của ngưỡng $K$ đó?
	\shortans[oly]{15}
	\loigiai{Khi thời gian $t$ kéo dài thì ta có 
		\[\lim\limits_{t\to +\infty}S(t)=\lim\limits_{t\to +\infty}\dfrac{10 t+5}{t+1}=10.\]
		Do đó $\lim\limits_{t\to +\infty}V(t)=\lim\limits_{S\to 10}V(S)=\lim\limits_{S\to 10}\dfrac{5 S}{S+2}=\dfrac{5\cdot10}{10+2}=\dfrac{25}{6}.$	 Suy ra $K=\dfrac{25}{6}$.	\\
		Ta có $V(S)=90\%\cdot\dfrac{25}{6}=\dfrac{15}{4}$.\\
		Khi đó $\dfrac{5 S}{S+2}=\dfrac{15}{4}\Leftrightarrow S=6$.\\
		Suy ra $S=6\Leftrightarrow\dfrac{10 t+5}{t+1}=6\Leftrightarrow 10t+5=6t+6\Leftrightarrow t=\dfrac{1}{4}=0{,}25$ h $=15$ phút.
	}
\end{ex}

\begin{ex}%[2H5V1-3]%[Dự án đề TT Toán NH25-26 - Đợt 4 -Lê Hoàng Hạc ]%[Sở Giáo Dục Hưng Yên]%[GVPB: Lê Minh]
	Trong không gian với hệ trục tọa độ $Oxyz$, cho mặt phẳng $(P)$ đi qua điểm $M(2;4;8)$ và cắt các tia $Ox$, $Oy$, $Oz$ lần lượt tại $A(a;0;0)$, $B(0;b;0)$, $C(0;0;c)$ sao cho $OA = 2OB = 4OC$. Tính $T = a+b+c$.
	\shortans{73{,}5}
	\loigiai{
		Vì $A$, $B$, $C$ thuộc các tia $Ox$, $Oy$, $Oz$ nên $a > 0$, $b > 0$, $c > 0$. \\
		Theo giả thiết $a = 2b = 4c \Rightarrow b = \dfrac{a}{2}$, $c = \dfrac{a}{4}$.\\
		Phương trình mặt phẳng $(P)$ theo đoạn chắn là $\dfrac{x}{a} + \dfrac{y}{b} + \dfrac{z}{c} = 1 \Leftrightarrow \dfrac{x}{a} + \dfrac{2y}{a} + \dfrac{4z}{a} = 1$.\\
		Thay tọa độ $M(2;4;8)$ vào phương trình trên ta được $\dfrac{2}{a} + \dfrac{8}{a} + \dfrac{32}{a} = 1 \Rightarrow a = 42$.\\
		Suy ra $b = 21$, $c = 10{,}5$.\\
		Vậy $T = a + b + c = 42 + 21 + 10{,}5 = 73{,}5$.
	}
\end{ex}

\begin{ex}%[0D0C2-7]%[Dự án đề TT Toán NH25-26 - Đợt 4 -Lê Hoàng Hạc ]%[Sở Giáo Dục Hưng Yên]%[GVPB: Lê Minh]
	Cho tập hợp $X = \{1;2;3;4;5;6;7;8\}$. Gọi $S$ là tập hợp tất cả các số tự nhiên có $4$ chữ số được lập từ các chữ số thuộc tập $X$. Chọn ngẫu nhiên một số từ tập hợp $S$. Xác suất để chọn được một số chia hết cho $3$ bằng $\dfrac{a}{b}$ (với $a,b \in \mathbb{N}^*$, $\dfrac{a}{b}$ là phân số tối giản). Tính $T=a+b$.
	\shortans{2731}
	\loigiai{
		Số phần tử của tập hợp $S$ là $n(\Omega) = 8^4 = 4096$.\\
		Chia các phần tử của tập $X$ thành $3$ nhóm dựa trên số dư khi chia cho $3$:
		\begin{itemize}
			\item Nhóm chia hết cho $3$: $X_0 = \{3; 6\}$ gồm $2$ phần tử.
			\item Nhóm chia $3$ dư $1$: $X_1 = \{1; 4; 7\}$ gồm $3$ phần tử.
			\item Nhóm chia $3$ dư $2$: $X_2 = \{2; 5; 8\}$ gồm $3$ phần tử.
		\end{itemize}
		Gọi $A$ là biến cố \lq\lq Số được chọn chia hết cho $3$\rq\rq.\\
		Một số chia hết cho $3$ khi và chỉ khi tổng các chữ số của nó chia hết cho $3$.\\
		Giả sử số được lập lấy $x$ chữ số từ $X_0$, $y$ chữ số từ $X_1$ và $z$ chữ số từ $X_2$.\\
		Ta có điều kiện: $\heva{&x + y + z = 4 \\ &(y + 2z) \,\vdots\, 3}$ (với $x, y, z \in \mathbb{N}$ và $0 \le x, y, z \le 4$).\\
		Ta có các trường hợp thỏa mãn như sau
		\begin{itemize}
			\item \textbf{TH1:} $x = 4, y = 0, z = 0$ ($4$ chữ số thuộc $X_0$).
			Số cách lập là $2^4 = 16$ cách.
			\item \textbf{TH2:} $x = 1, y = 3, z = 0$ ($1$ chữ số thuộc $X_0$, $3$ chữ số thuộc $X_1$).
			Số cách lập là $\mathrm{C}_4^1 \times 2^1 \times 3^3 = 216$ cách.
			\item \textbf{TH3:} $x = 1, y = 0, z = 3$ ($1$ chữ số thuộc $X_0$, $3$ chữ số thuộc $X_2$).
			Số cách lập là $\mathrm{C}_4^1 \times 2^1 \times 3^3 = 216$ cách.
			\item \textbf{TH4:} $x = 2, y = 1, z = 1$ ($2$ chữ số thuộc $X_0$, $1$ thuộc $X_1$, $1$ thuộc $X_2$).
			Số cách lập là $\mathrm{C}_4^2 \times \mathrm{C}_2^1 \times 2^2 \times 3^1 \times 3^1 = 432$ cách.
			\item \textbf{TH5:} $x = 0, y = 2, z = 2$ ($2$ chữ số thuộc $X_1$, $2$ chữ số thuộc $X_2$).
			Số cách lập là $\mathrm{C}_4^2 \times 3^2 \times 3^2 = 486$ cách.
		\end{itemize}
		Từ $5$ trường hợp trên, số các số chia hết cho $3$ trong tập $S$ là $$n(A) = 16 + 216 + 216 + 432 + 486 = 1366.$$
		Xác suất của biến cố $A$ là $\mathrm{P}(A) = \dfrac{n(A)}{n(\Omega)} = \dfrac{1366}{4096} = \dfrac{683}{2048}$.\\
		Vì $\dfrac{683}{2048}$ là phân số tối giản nên $a = 683, b = 2048$.\\
		Vậy $T = a + b = 683 + 2048 = 2731$.
	}
\end{ex}
=======
\def\thoigian{90}%--Thời gian
\de{ĐỀ THI THỬ LẦN 1 NĂM HỌC 2025-2026}{SỞ GIÁO DỤC HƯNG YÊN}



\begin{center}
	\textbf{PHẦN 1 - Câu trắc nghiệm nhiều phương án lựa chọn.}
\end{center}
\setcounter{ex}{0}

\Opensolutionfile{ans}[ans-ABCD]




\begin{ex}%[Đề thi thử môn Toán Sở GD Hưng Yên - lần 1, năm học 2025-2026]%[GV biên soạn: Nhật Thiện - GV phản biện: Lê Minh]%[2H5N1-2]
	Trong không gian với hệ tọa độ $Oxyz$, cho mặt phẳng $(P)$ có phương trình $x+2y-z-2=0$. Vectơ nào sau đây là một vectơ pháp tuyến của mặt phẳng $(P)$?
	\choice
	{\True $\overrightarrow{c}(1;2;-1)$}
	{$\overrightarrow{b}(1;-2;-1)$}
	{$\overrightarrow{d}(1;-2;1)$}
	{$\overrightarrow{a}(1;2;1)$}
	\loigiai{
		Một vectơ pháp tuyến của mặt phẳng $(P)$ là $\overrightarrow{c}(1;2;-1)$.
	}
\end{ex}
\begin{ex}%[Đề thi thử môn Toán Sở GD Hưng Yên - lần 1, năm học 2025-2026]%[GV biên soạn: Nhật Thiện - GV phản biện: Lê Minh]%[1D6H2-1]
	Cho $\log_ab=5$. Giá trị của biểu thức $\log_a\left(a^2b\right)$ bằng
	\choice
	{$20$}
	{$25$}
	{$10$}
	{\True $7$}
	\loigiai{
		Ta có $\log_a\left(a^2b\right)=\log_a a^2+\log_a b=2+\log_ab=2+5=7$.
	}
\end{ex}
\begin{ex}%[Đề thi thử môn Toán Sở GD Hưng Yên - lần 1, năm học 2025-2026]%[GV biên soạn: Nhật Thiện - GV phản biện: Lê Minh]%[2H2H2-4]
	Trong không gian với hệ tọa độ $Oxyz$, khoảng cách giữa hai điểm $A(-2;-3;9)$ và $B(-2;3;1)$ bằng
	\choice
	{$8$}
	{$6$}
	{$100$}
	{\True $10$}
	\loigiai{
		Khoảng cách $AB=\sqrt{(-2+2)^2+(3+3)^2+(1-9)^2}=10$.
	}
\end{ex}
\begin{ex}%[Đề thi thử môn Toán Sở GD Hưng Yên - lần 1, năm học 2025-2026]%[GV biên soạn: Nhật Thiện - GV phản biện: Lê Minh]%[2H5H1-3]
	Trong không gian với hệ tọa độ $Oxyz$, phương trình mặt phẳng $(P)$ đi qua điểm $I(2;-1;3)$ và có véctơ pháp tuyến $\overrightarrow{n}(2;-2;1)$ là
	\choice
	{$2x-y+3z-9=0$}
	{$2x-2y+z+9=0$}
	{$2x-y+3z+9=0$}
	{\True $2x-2y+z-9=0$}
	\loigiai{
		Phương trình mặt phẳng $(P)$ đi qua điểm $I(2;-1;3)$ và có véctơ pháp tuyến $\overrightarrow{n}(2;-2;1)$ là
		$$2(x-2)-2(y+1)+1(z-3)=0\Leftrightarrow 2x-2y+z-9=0.$$
	}
\end{ex}
\begin{ex}%[Đề thi thử môn Toán Sở GD Hưng Yên - lần 1, năm học 2025-2026]%[GV biên soạn: Nhật Thiện - GV phản biện: Lê Minh]%[2D1N2-2]
	\immini[thm]{Cho hàm số $y=ax^3+bx^2+cx+d$ ($a\ne 0$) có đồ thị như hình vẽ bên. Điểm cực tiểu của đồ thị hàm số đã cho là
		\choice
		{$y=3$}
		{$y=-1$}
		{\True $x=1$}
		{$x=-1$}}{\begin{tikzpicture}[scale=1, font=\footnotesize, line join=round, line cap=round, >=stealth]
			\begin{scope}[scale=.8]
				\draw[->] (-2.5, 0) -- (2.5, 0) node[below] {$x$};
				\draw[->] (0, -2) -- (0, 4.2) node[left] {$y$};
				\node[below left] at (0, 0) {$O$};
				\begin{scope}
					\clip (-2,-2) rectangle (2,4);
					\draw[domain=-2.15:2.15] plot (\x, {(\x)^3 - 3*(\x) + 1});
					\draw[dashed] (-1, 0) node[below] {$-1$} -- (-1, 3) -- (0, 3) node[right] {$3$};
					\draw[dashed] (1, 0) node[above] {$1$} -- (1, -1) -- (0, -1) node[left] {$-1$};
				\end{scope}
				\node[right] at (0.05, 1) {$1$};
			\end{scope}
	\end{tikzpicture}}
	\loigiai{
		Dựa vào đồ thị, điểm cực tiểu của hàm số đã cho là $x=1$.
	}
\end{ex}
\begin{ex}%[Đề thi thử môn Toán Sở GD Hưng Yên - lần 1, năm học 2025-2026]%[GV biên soạn: Nhật Thiện - GV phản biện: Lê Minh]%[1H8H6-2]
	\immini[thm]{Cho hình chóp $S.ABC$ có đáy $ABC$ là tam giác vuông tại $B$, $SA\perp (ABC)$, $SA=AB=a$ (hình vẽ minh họa). Tính số đo của góc nhị diện $[S,BC,A]$.
		\choice
		{$60^\circ$}
		{\True $45^\circ$}
		{$30^\circ$}
		{$90^\circ$}}{
		\begin{tikzpicture}[scale=1, font=\footnotesize, line join=round, line cap=round, >=stealth]
			\path (0,0) coordinate (A)	(1,-1) coordinate (B) (4,0) coordinate (C) (A)+(0,3) coordinate (S);
			\draw (S)--(A)--(B)--(C)--cycle (S)--(B);
			\draw[dashed] (A)--(C);
			\foreach \x/\y in {A/180,B/-90,C/0,S/90} \draw[fill=black] (\x) circle (1pt) node[shift=(\y:0.3cm)] {$\x$};
		\end{tikzpicture}
	}
	\loigiai{
		Ta có $BC\perp AB$ và $BC\perp SA$ suy ra $BC\perp SB$.\\
		Mặt khác $AB\perp BC$.\\
		Do đó, $\widehat{SBA}$ là góc phẳng nhị diện $[S,BC,A]$.\\
		Xét tam giác $SAB$ là tam giác vuông cân tại $A$ nên $\widehat{SBA}=45^\circ$.
		Vậy số đo của góc nhị diện $[S,BC,A]$ bằng $45^\circ$.
	}
\end{ex}
%%Câu 7%%
\begin{ex}%[2D1N4-1]%[Dự án C đề thi thử TN Môn Toán NH25-26 - Dot 4 - Hoàng Tuấn]%[Sở Hưng Yên]%[GVPB - Tên GV]		
	Đường tiệm cận đứng của đồ thị hàm số $y = \dfrac{2x + 3}{x - 2}$ có phương trình là
	\choice
	{$y = 2$}
	{\True $x = 2$}
	{$y = -\dfrac{3}{2}$}
	{$x = -\dfrac{3}{2}$}
	\loigiai{
		Đường tiệm cận đứng của đồ thị hàm số $y = \dfrac{2x + 3}{x - 2}$ có phương trình là $x=2$.
	}
\end{ex}

%%Câu 8%%
\begin{ex}%[2D1N1-1]%[Dự án C đề thi thử TN Môn Toán NH25-26 - Dot 4 - Hoàng Tuấn]%[Sở Hưng Yên]%[GVPB - Tên GV]
	Hàm số $y = x^3 - 3x + 2026$ nghịch biến trên khoảng nào dưới đây?
	\choice
	{$(1; +\infty)$}
	{$(-\infty; 1)$}
	{$(-2; 2)$}
	{\True $(-1; 1)$}
	\loigiai{Tập xác định $\mathscr{D}=\mathbb{R}$.\\
		Ta có $y' = 3x^2 - 3$.\\
		$y'=0\Leftrightarrow \hoac{&x=1\\&x=-1.}$\\
		Bảng biến thiên
		\begin{center}
			\begin{tikzpicture}
				\tkzTabInit[nocadre=false,lgt=1.2,espcl=2.5,deltacl=0.6]
				{$x$/0.6, $y’$/0.6, $y$/2}%
				{$-\infty$,$-1$,$1$,$+\infty$}
				\tkzTabLine{,+,0,-,0,+,}
				\tkzTabVar{-/$-\infty$, +/$2028$,-/$2024$,+/$+\infty$}
			\end{tikzpicture}
		\end{center}
		Từ bảng biến thiên ta thấy  hàm số nghịch biến trên khoảng $(-1; 1)$.
	}
\end{ex}

%%Câu 9%%
\begin{ex}%[2D4H2-1]%[Dự án C đề thi thử TN Môn Toán NH25-26 - Dot 4 - Hoàng Tuấn]%[Sở Hưng Yên]%[GVPB - Tên GV]
	Cho hàm số $y = f(x)$ liên tục trên $\mathbb{R}$, thỏa mãn $\displaystyle\int_{17}^{2026} f(x)\, \mathrm{\,d}x = 8$ và $\displaystyle\int_{89}^{2026} f(x)\, \mathrm{\,d}x = 10$. Giá trị của biểu thức $\displaystyle\int_{17}^{89} f(x)\, \mathrm{\,d}x$ bằng
	\choice
	{$18$}
	{$2026$}
	{$2$}
	{\True $-2$}
	\loigiai{
		Ta có
		\begin{eqnarray*}
			\displaystyle\int_{17}^{89} f(x) \mathrm{\,d}x &=& \displaystyle\int_{17}^{2026} f(x) \mathrm{\,d}x + \displaystyle\int_{2026}^{89} f(x) \mathrm{\,d}x \\
			&=& \displaystyle\int_{17}^{2026} f(x) \mathrm{\,d}x - \displaystyle\int_{89}^{2026} f(x) \mathrm{\,d}x\\
			&=&8-10\\
			&=&-2.
		\end{eqnarray*}
		
	}
\end{ex}

%%Câu 10%%
\begin{ex}%[2D4N1-3]%[Dự án C đề thi thử TN Môn Toán NH25-26 - Dot 4 - Hoàng Tuấn]%[Sở Hưng Yên]%[GVPB - Tên GV]
	Họ nguyên hàm của hàm số $f(x) = \sin x$ là
	\choice
	{$-\sin x + C$}
	{\True $-\cos x + C$}
	{$\sin x + C$}
	{$\cos x + C$}
	\loigiai{
		Họ nguyên hàm của hàm số $f(x)=\sin x$ là $$\displaystyle\int \sin x \,\mathrm{\,d}x = -\cos x + C.$$
	}
\end{ex}


%%Câu 11%%
\begin{ex}%[1D5H1-3]%[Dự án C đề thi thử TN Môn Toán NH25-26 - Dot 4 - Hoàng Tuấn]%[Sở Hưng Yên]%[GVPB - Tên GV]
	Kết quả kiểm tra môn Toán của học sinh lớp $12\text{B}$ được thống kê như sau:
	\begin{center}
		\begin{tabular}{|l|c|c|c|c|c|}
			\hline
			Điểm & $[0; 2)$ & $[2; 4)$ & $[4; 6)$ & $[6; 8)$ & $[8; 10)$ \\
			\hline
			Số học sinh & $1$ & $4$ & $15$ & $16$ & $4$ \\
			\hline
		\end{tabular}
	\end{center}
	Số trung bình cộng của mẫu số liệu ghép nhóm trên bằng
	\choice
	{$6{,}1$}
	{$6$}
	{\True $5{,}9$}
	{$5{,}8$}
	\loigiai{
		Giá trị đại diện của các nhóm lần lượt là $c_1 = 1$, $c_2 = 3$, $c_3 = 5$, $c_4 = 7$, $c_5 = 9$.\\
		Tổng số học sinh $n= 1 + 4 + 15 + 16 + 4 = 40$.\\
		Số trung bình cộng là:
		$\overline{x} = \dfrac{1 \cdot 1 + 3 \cdot 4 + 5 \cdot 15 + 7 \cdot 16 + 9 \cdot 4}{40} = \dfrac{236}{40} = 5{,}9$.
	}
\end{ex}

%%Câu 12%%
\begin{ex}%[1D5H2-3]%[Dự án C đề thi thử TN Môn Toán NH25-26 - Dot 4 - Hoàng Tuấn]%[Sở Hưng Yên]%[GVPB - Tên GV]
	Bảng dưới đây biểu diễn mẫu số liệu ghép nhóm về chiều cao (đơn vị: cm) của $40$ cây cà chua trong vườn ươm:
	\begin{center}
		\begin{tabular}{|c|c|c|c|c|}
			\hline
			Nhóm & $[20; 25)$ & $[25; 30)$ & $[30; 35)$ & $[35; 40)$ \\
			\hline
			Tần số & $5$ & $9$ & $25$ & $1$ \\
			\hline
		\end{tabular}
	\end{center}
	Tứ phân vị thứ ba $Q_3$ (đơn vị: cm) của mẫu số liệu ghép nhóm trên bằng
	\choice
	{$33$}
	{\True $33{,}2$}
	{$33{,}1$}
	{$33{,}4$}
	\loigiai{
		Tổng số mẫu $n = 40$.\\
		Do  $\dfrac{3n}{4} =\dfrac{3\cdot 40}{4}= 30$.\\
		Vị trí của tứ phân vị thứ ba là của $Q_3=\dfrac{x_{30}+x_{31}}{2}$.\\
		Giá trị này rơi vào nhóm $[30;35)$. Đây là nhóm chứa $Q_3$.
		Áp dụng công thức ta có
		\begin{eqnarray*}
			Q_3 &=&u_m + \dfrac{\dfrac{3n}{4} - C}{n_m} \cdot (u_{m+1} - u_m)\\
			&=&30 + \dfrac{30 - 14}{25} \cdot 5\\
			&=&30 + \dfrac{16}{5}\\
			&=&30 + 3{,}2\\
			&=&33{,}2.
		\end{eqnarray*}
	}
\end{ex}



\Closesolutionfile{ans}


\begin{center}
	\textbf{PHẦN 2 - Câu trắc nghiệm đúng sai. Trong mỗi ý a,b,c,d ở mỗi câu, thí sinh chọn đúng hoặc sai}
\end{center}
\setcounter{ex}{0}
\Opensolutionfile{ans}[ans-DS]
\begin{ex}%[Đề thi thử môn Toán SGD Hưng Yên Lần 1 năm học 2025-2026]%[Xuan Vy Pham - Lê Minh]%[2D4V3-3]
	Cho hàm số $f(x)= \sin x +2x$. Hàm số $F(x)$ là một nguyên hàm của hàm số $f(x)$ trên $\mathbb{R}$ và thỏa mãn $F(0)=4$.
	\choiceTF
	{\True Tiếp tuyến của đồ thị hàm số $y=F(x)$ tại điểm có hoành độ $x = \pi$ có hệ số góc $k=2\pi$}
	{\True Họ nguyên hàm của hàm số $f(x)$ là $-\cos x+x^2+C$ (với $C$ là hằng số)}
	{$F(x)= -\cos x+x^2+4$}
	{\True Thể tích khối tròn xoay sinh bởi miền phẳng $D$ giới hạn bởi đồ thị hàm số $y=f(x)$, trục hoành và hai đường thẳng có phương trình $x=0$, $x= \dfrac{\pi}{2}$ bằng $31$ (đvtt) (\textit{không làm tròn kết quả của các phép toán trung gian, chỉ làm tròn kết quả phép toán cuối cùng đến hàng đơn vị})}
	\loigiai{
		\begin{itemchoice}
			\itemch  
			Vì $F(x)$ là một nguyên hàm của hàm số $f(x)$ trên $\mathbb{R}$ nên $F'(x)=f(x)$, $\forall x \in \mathbb{R}$.\\
			Tiếp tuyến của đồ thị hàm số $y=F(x)$ tại điểm có hoành độ $x = \pi$ có hệ số góc là $$k=F'(\pi)=f(\pi)=\sin \pi +2 \pi = 2\pi.$$
			\itemch  Họ nguyên hàm của hàm số $f(x)$ là 
			$$F(x)=\displaystyle \int f(x) \mathrm{\,d}x= \displaystyle \int \left(\sin x +2x\right)\mathrm{\,d}x=-\cos x +x^2+C.$$
			\itemch Vì $F(0)=4$ nên $-\cos 0+0^2+C=4 \Rightarrow C=5$.\\
			Vậy $F(x)=-\cos x +x^2+5$.
			\itemch Thể tích khối tròn xoay sinh bởi miền phẳng $D$ giới hạn bởi đồ thị hàm số $y=f(x)$, trục hoành và hai đường thẳng có phương trình $x=0$, $x= \dfrac{\pi}{2}$ là
			$$V=\pi \displaystyle \int \limits_{0}^{\tfrac{\pi}{2}} \left[f(x)\right]^2 \mathrm{\,d}x = \pi \displaystyle \int \limits_{0}^{\tfrac{\pi}{2}} \left( \sin x+2x\right)^2 \mathrm{\,d}x  \approx 31\; \text{(đvtt)}.$$
		\end{itemchoice}
	}
\end{ex}
\begin{ex}%[Đề thi thử môn Toán SGD Hưng Yên Lần 1 năm học 2025-2026]%[Xuan Vy Pham - Lê Minh]%[2D1V5-8]	
	Một hộ gia đình muốn xây dựng một bể chứa nước có dạng hình hộp chữ nhật không nắp, thể tích $V=12$m$^3$. Đáy bể là hình chữ nhật có chiều dài gấp đôi chiều rộng. Giá thuê nhân công và vật liệu xây đáy là $500$ nghìn đồng/m$^2$, xây thành bể là $300$ nghìn đồng/m$^2$. Gọi $x$ là chiều rộng của đáy bể, $h$ là chiều cao của bể ($x>0$, $h>0$, đơn vị: m).
	\choiceTF
	{\True Thể tích của bể nước được tính bằng công thức $V=2x^2h$ (m$^3$)}
	{\True Chiều cao của bể nước tính theo $x$ là $h=\dfrac{6}{x^2}$ (m)}
	{Tổng chi phí xây dựng bể nước là $T(x)=500x^2+\dfrac{10\,800}{x}$ (nghìn đồng)}
	{Tổng chi phí tối thiểu để xây dựng bể là $9\,324$ (nghìn đồng) (\textit{không làm tròn kết quả của các phép toán trung gian, chỉ làm tròn kết quả phép toán cuối cùng đến hàng đơn vị})}
	\loigiai{\begin{center}
			\begin{tikzpicture}[line join=round, line cap=round,>=stealth,font=\footnotesize,scale=1]
				\def\a{3} 
				\def\b{1.8}
				\def\h{3.5}
				\path 	(0:0) coordinate (A)
				++(0:\a) coordinate (D)
				++(-130:\b) coordinate (C)
				($(A)+(C)-(D)$) coordinate (B)
				($(A)+(90:\h)$) coordinate (A')
				($(B)+(90:\h)$) coordinate (B')
				($(C)+(90:\h)$) coordinate (C')
				($(D)+(90:\h)$) coordinate (D');
				\draw[dashed,thick] 	(B)--(A)--(D)	(A)--(A');
				\draw[thick] (C)--(C') 	(D)--(D') 	(B)--(B') 	(B)--(C)--(D) (A')--(B')--(C')--(D')--cycle;
				\foreach \x/\g in {A/180,B/180,C/0,D/0,A'/180,B'/180,C'/0,D'/0}
				\fill[black] 	(\x) circle (1pt)
				($(\g:4mm)+(\x)$) node {$\x$};	
				\draw ($(A)!1/2!(D)$) node[above]{$2x$};
				\draw ($(C)!1/2!(D)$) node[right]{$x$};
				\draw ($(A)!1/2!(A')$) node[right]{$h$};
			\end{tikzpicture}
		\end{center}
		\begin{itemchoice}
			\itemch Thể tích của bể nước được tính bằng công thức $V=x \cdot (2x) \cdot h =2x^2h$ (m$^3$).
			\itemch Ta có $V=12 \Rightarrow 2x^2 h=12 \Rightarrow h =\dfrac{6}{x^2}$ (m).
			\itemch Chi phí để xây dựng đáy bể là $500 \cdot x \cdot 2x  = 1\, 000x^2$ (nghìn đồng).\\
			Chi phí để xây dựng thành bể là $$300  \left(2x  h+ 4x  h\right)=1\,800 xh= 1\,800 x\cdot \dfrac{6}{x^2} = \dfrac{19 \, 800}{x^2} \; \text{(nghìn đồng)}.$$
			Tổng chi phí xây dựng bể nước là $T(x)=1\,000x^2+\dfrac{10\,800}{x}$ (nghìn đồng).
			\itemch
			Ta có $T'(x)=2 \, 000x-\dfrac{10 \, 800}{x^2}$.\\
			$T'(x) = 0 \Leftrightarrow 2 \, 000x-\dfrac{10 \, 800}{x^2} = 0 \Leftrightarrow 2\,000x^3 - {10 \, 800}=0 \Leftrightarrow x = \dfrac{3}{\sqrt[3]{5}}$. 
			\begin{center}
				\begin{tikzpicture}
					\tkzTabInit[nocadre,lgt=1.2,espcl=2.5,deltacl=0.6]
					{$x$ /1.4, $T’(x)$ /0.6, $T(x)$ /2}
					{$0$,$\dfrac{3}{\sqrt[3]{5}}$,$+\infty$}
					\tkzTabLine{ ,-,z,+,}
					\tkzTabVar{+/$+\infty$,-/$T\left(\dfrac{3}{\sqrt[3]{5}}\right)$,+/$+\infty$}
				\end{tikzpicture}
			\end{center}
			Vì $\min \limits_{(0,+\infty)} T(x)=T\left(\dfrac{3}{\sqrt[3]{5}}\right)$ nên
			tổng chi phí tối thiểu để xây dựng bể là $$T \left(\dfrac{3}{\sqrt[3]{5}}\right) \approx 9\,234 \; \text{(nghìn đồng)}.$$
		\end{itemchoice}
	}
\end{ex}


\begin{ex}%[2H2N2-3]%[2H2H2-2]%[2H2H2-4]%[2H5H1-3]%[Dự án C  đề thi thử TN Môn Toán NH25-26 - Đợt 4- Phan Trung Hiếu]%[TT-SGD-Hưng Yên L1 - NH25-26]%[GVPB: Nguyễn Hoài Nam]
	Trong không gian với hệ tọa độ $Oxyz$, cho hình hộp $ABCD.A'B'C'D'$ có tọa độ các đỉnh $A(1;-1;3)$, $B(0;2;4)$, $D(2;-1;1)$ và $A'(0;1;2)$.
	\choiceTF
	{\True Tọa độ của vectơ $\overrightarrow{AD} = (1; 0; -2)$}
	{\True Tọa độ đỉnh $B'(-1; 4; 3)$}
	{Góc giữa hai vectơ $\overrightarrow{AB}$ và $\overrightarrow{AD}$ là góc nhọn}
	{Phương trình mặt phẳng $(CB'D')$ có dạng $ax + by + cz - 7 = 0$. Khi đó: $a + b - c = 10$}
	\loigiai{
		\begin{center}
			\begin{tikzpicture}[font=\footnotesize, thick, >=latex]
				\path 
				(0,0) coordinate (A)
				(4,0) coordinate (D)
				(-1,-1) coordinate (B)
				($(B)+(D)-(A)$) coordinate (C)
				(1,3) coordinate (A')
				($(B)+(A')-(A)$) coordinate (B')
				($(D)+(A')-(A)$) coordinate (D')
				($(C)+(A')-(A)$) coordinate (C')
				;
				\draw (A')--(B')--(C')--(D')--cycle;
				\draw (B)--(C)--(D) (B)--(B') (C)--(C') (D)--(D');
				\draw[dashed] (B)--(A)--(D) (A)--(A');
				\foreach \x/\g in {A/160,B/180,C/-45,D/0, A'/90,B'/180,C'/90,D'/0}
				\fill[black] 	(\x) circle (1pt)
				($(\g:3mm)+(\x)$) node {$\x$};
			\end{tikzpicture}
		\end{center}
		\begin{itemchoice}
			\itemch $\overrightarrow{AD} = (1;0;-2)$.
			\itemch Ta có $\overrightarrow{AB} = \overrightarrow{A'B'} \Leftrightarrow B' = A' + B - A = \left(-1;4;3\right)$. Vậy $B'\left(-1;4;3\right)$.
			\itemch $\overrightarrow{AB} = (-1;3;1)$. Ta có
			\begin{equation*}
				\cos\left(\overrightarrow{AB},\overrightarrow{AD}\right) = \dfrac{\overrightarrow{AB}\cdot\overrightarrow{AD}}{\left|\overrightarrow{AB}\right|.\left|\overrightarrow{AD}\right|}=\dfrac{-1\cdot1+3\cdot0+1\cdot(-2)}{\sqrt{1^2+0^2+(-2)^2}\cdot\sqrt{(-1)^2+3^2+1^2}} = \dfrac{-3}{\sqrt{55}}<0.
			\end{equation*}
			Do đó góc giữa hai vectơ $\overrightarrow{AB}$ và $\overrightarrow{AD}$ là góc tù.
			\itemch $\overrightarrow{BC} = \overrightarrow{AD}\Leftrightarrow C = B + D- A = (1;2;2)$. Vậy $C(1;2;2)$.\\
			Tương tự, $\overrightarrow{AD} = \overrightarrow{A'D'}\Leftrightarrow D'(1;1;0)$.\\
			Ta có $\overrightarrow{CB'}=(-2;2;1)$, $\overrightarrow{CD'}=(0;-1;-2)$.\\
			Véc-tơ pháp tuyến của mặt phẳng $(CB'D')$ là $\overrightarrow{n}=\left[\overrightarrow{CB'},\overrightarrow{CD'}\right] = (-3;-4;2)$.\\
			Khi đó, phương trình mặt phẳng $(CB'D')$ có dạng $-3x-4y+2z+D=0$.\\
			Vì $C\in(C'B'D')$ nên $-3-8+4+D=0\Leftrightarrow D=7$.\\
			Vậy $(CB'D')\colon 3x+4y-2z-7=0$. Suy ra $a=3$, $b=4$, $c=-2$.\\
			Khi đó, $a+b-c=9$.
		\end{itemchoice}
	}
\end{ex}

\begin{ex}%[0D2V2-3]%[Dự án C  đề thi thử TN Môn Toán NH25-26 - Đợt 4- Phan Trung Hiếu]%[TT-SGD-Hưng Yên L1 - NH25-26]%[GVPB: Nguyễn Hoài Nam]
	Một xưởng sơn dự định sản xuất hai loại sơn là sơn nội thất và sơn ngoài trời. Nguyên liệu để sản xuất gồm hai loại A và B với trữ lượng lần lượt là $6$ tấn và $8$ tấn. Để sản xuất một tấn sơn nội thất cần $2$ tấn nguyên liệu A và $1$ tấn nguyên liệu B. Để sản xuất một tấn sơn ngoài trời cần $1$ tấn nguyên liệu A và $2$ tấn nguyên liệu B. Kết quả nghiên cứu thị trường cho thấy nhu cầu sơn nội thất không quá $2$ tấn và không nhiều hơn nhu cầu sơn ngoài trời $1$ tấn. Giá bán một tấn sơn nội thất là $60$ triệu đồng, một tấn sơn ngoài trời là $30$ triệu đồng (giả sử rằng số sơn sản xuất ra đều được bán hết).
	\choiceTF
	{Gọi $x$; $y$ lần lượt là số tấn sơn nội thất và sơn ngoài trời cần sản xuất $x \ge 0$, $y \ge 0$}
	{Biểu thức doanh thu $F(x, y) = 30x + 60y$ (triệu đồng)}
	{\True Doanh thu lớn nhất của xưởng là $180$ triệu đồng}
	{Trong các phương án sản xuất đem lại doanh thu lớn nhất, biết rằng tổng số lượng sơn cả hai loại dự định sản xuất không quá $4{,}5$ tấn. Khi đó lượng sơn nội thất cần sản xuất ít nhất là $1{,}4$ tấn}
	
	\loigiai{
		\begin{itemchoice}
			\itemch Gọi $x$; $y$ lần lượt là số tấn sơn nội thất và sơn ngoài trời cần sản xuất.\\
			Điều kiện: $0\leq x \leq 2$, $y\geq 0$ và $x\leq y+1$.
			\immini[thm]{
				Ta có hệ phương trình $\heva{&0\leq x\leq 2\\&y \leq 0\\&x-y\leq1\\&2x+y\leq 6&\\&x+2y\leq 8\\&x+y\leq4{,}5.}$\\
				Miền nghiệm là miền đa giác $OAFBCDE$.
				\begin{itemize}
					\item Tại $A(0,4)$ ta có $F(0,4) = 120$.
					\item Tại $F\left(\dfrac{3}{2};3\right)$ ta có $F\left(\dfrac{3}{2};3\right) = 180$.
					\item Tại $C(2,2)$, ta có $F(2,2) = 180$.
					\item Tại $D(2,1)$ ta có $F(2,1) = 150$.
					\item Tại $E(1,0)$ ta có $F(1,0) = 60$.
				\end{itemize}
			}{
				\begin{tikzpicture}[font=\footnotesize, line cap=round, line join = round, thick, >=stealth, scale=.7,transform shape]
					\def\xmin{-1}
					\def\xmax{9}
					\def\ymin{-1}
					\def\ymax{9}
					\clip (\xmin,\ymin-.5) rectangle (\xmax,\ymax);
					\draw[->] (\xmin,0)--(0,0) node[below left]{$O$} -- (\xmax,0)node[above left]{$x$};
					\draw[->] (0,\ymin-.5)--(0,\ymax)node[below right]{$y$};
					\draw (2,\ymin-.5)--(2,\ymax);
					\draw[domain = \xmin:\xmax] plot (\x, {\x -1});
					\draw[domain = \xmin:\xmax] plot (\x, {-2*\x +6});
					\draw[domain = \xmin:\xmax] plot (\x, {-\x/2 +4});
					\draw[domain = \xmin:\xmax] plot (\x, {-\x +4.5});
					\fill[pattern= north west lines, pattern color =  cyan] (\xmin,0) rectangle (\xmax,\ymin-.5);
					\fill[pattern= north west lines, pattern color =  red] (\xmin,\ymin-.5) rectangle (0,\ymax);
					\fill[pattern= north east lines, pattern color =  blue] (2,\ymin-.5) rectangle (\xmax,\ymax);
					\fill[pattern= north east lines, pattern color =  yellow] (2,\ymin-.5) rectangle (\xmax,\ymax);
					\fill[pattern= north west lines, pattern color =  pink] (\xmin,\ymax) -- plot[domain = \xmin:\xmax] (\x, {-\x/2 +4}) -- (\xmax,\ymax);
					%				\fill[pattern= north east lines, pattern color =  magenta!30] (\xmin,\ymax) -- plot [domain = \xmin:\xmax](\x, {-2*\x +6}) -- (\xmax,\ymax);
					\fill[pattern= north east lines, pattern color =  magenta!30] (\xmin,\ymax) -- plot [domain = \xmin:\xmax](\x, {-\x +4.5}) -- (\xmax,\ymax);
					\foreach \x/\y/\pos/\name in {0/4/200/A,1.5/3/-120/B,2/2/0/C,2/1/-45/D,1/0/-90/E, 1/3.5/-120/F}
					\fill[black] 	(\x,\y) circle (1.5pt)
					($(\pos:3mm)+(\x,\y)$) node {$\name$};
				\end{tikzpicture}
			}
			\itemch Biểu thức doanh thu $F(x;y)=60x+30y$ (triệu đồng).
			%			\itemch Các phương án sản xuất đem lại doanh thu lớn nhất biết rằng tổng số lượng sơn cả hai loại dự định sản xuất không quá $4{,}5$ tấn khi sản xuất $1{,}5$ tấn sơn nội thất và $3$ tấn sơn ngoài trời. Khi đó, lượng sơn nội thất cần sản xuất ít nhất là $1{,}5$ tấn.
			\itemch Doanh thu lớn nhất của xưởng là $180$ triệu đồng.
			\itemch Doanh thu đạt giá trị lớn nhất tại điểm $F\left(\dfrac{3}{2};3\right)$ hoặc $C(2,2)$, khi đó lượng sơn nội thất cần sản xuất ít nhất là $1{,}5$ tấn. 
		\end{itemchoice}
		
	}
\end{ex}


\Closesolutionfile{ans}


\begin{center}
	\textbf{PHẦN 3 - Câu trắc nghiệm trả lời ngắn}
\end{center}
\setcounter{ex}{0}

\Opensolutionfile{ans}[ans-KQ]
\begin{ex}%[2D1H2-7]%[Dự án đề TT Toán NH 2025-2026 Đợt 4-Lê Minh Thiện Anh]%[Sở GD Hưng Yên -Lần 1 - Hưng Yên]%[GVPB: Lê Minh]
	Trong trò chơi mô phỏng xây dựng công viên giải trí Roller Coaster Tycoon, một kỹ sư thiết kế một đoạn đường ray hình lượn sóng. Trong hệ trục tọa độ $Oxy$, đoạn đường ray này được mô hình hóa bởi đồ thị của hàm số: $y = x^3 - 6x^2 + 9x + 1$, với $0 \le x \le 4$. Biết hệ trục tọa độ được thiết lập sao cho trục $Ox$ nằm ngang và mỗi đơn vị độ dài trên các trục tọa độ tương ứng với $2$ mét trên thực tế. Để kiểm tra độ ổn định của cấu trúc, kỹ sư sử dụng thiết bị laser để đo khoảng cách trực tiếp giữa hai điểm chốt kỹ thuật đặt tại điểm cực đại $A$ và điểm cực tiểu $B$ của đồ thị hàm số trên. Hãy tính độ dài thực tế của khoảng cách giữa hai điểm $A$ và $B$ \textit{(làm tròn kết quả cuối cùng đến hàng phần trăm theo đơn vị mét)}.
	\shortans[oly]{8{,}94}
	\loigiai{
		Ta có $y' = 3x^2 - 12x + 9$.\\
		$y' = 0 \Leftrightarrow 3(x^2 - 4x + 3) = 0 \Leftrightarrow \hoac{ &x = 1 \Rightarrow y = 5 \\ &x = 3 \Rightarrow y = 1.}$\\
		Tọa độ hai điểm cực trị là $A(1; 5)$ và $B(3; 1)$.\\
		Khoảng cách giữa hai điểm $A$, $B$ trên mặt phẳng tọa độ là
		$$AB = \sqrt{(3-1)^2 + (1-5)^2} = \sqrt{2^2 + (-4)^2} = \sqrt{20} = 2\sqrt{5}.$$
		Vì mỗi đơn vị độ dài tương ứng với $2$ mét thực tế, nên độ dài thực tế là
		$$L = 2\sqrt{5} \cdot 2 = 4\sqrt{5} \approx 8{,}94 \,(\text{m}).$$
	}
\end{ex}

\begin{ex}%[1D6H3-5]%[Dự án đề TT Toán NH 2025-2026 Đợt 4-Lê Minh Thiện Anh]%[Sở GD Hưng Yên -Lần 1 - Hưng Yên]%[GVPB: Lê Minh]
	Chỉ số giá tiêu dùng CPI là chỉ số tính theo phần trăm phản ánh mức thay đổi của giá hàng hóa tiêu dùng theo thời gian. Giả sử tại một quốc gia, chỉ số CPI sau $n$ năm được dự báo theo công thức: $CPI_n = CPI_0(1+g)^n$, trong đó $CPI_0$ là chỉ số tại thời điểm bắt đầu dự báo và $g$ là tỉ lệ lạm phát trung bình hằng năm. Biết rằng vào tháng 1 năm 2020 chỉ số CPI của quốc gia này là $118{,}09$ và tỉ lệ lạm phát duy trì ổn định ở mức $g = 3{,}21\%$/năm. Hãy tính chỉ số CPI dự báo của quốc gia này vào tháng 1 năm $2030$ \textit{(làm tròn kết quả cuối cùng đến hàng đơn vị)}.
	\shortans[oly]{162}
	\loigiai{
		Thời gian từ tháng $1$ năm $2020$ đến tháng $1$ năm $2030$ là $n = 10$ năm.\\
		Áp dụng công thức với $CPI_0 = 118{,}09$, $g = 3,21\% = 0{,}0321$ và $n = 10$, ta có\\
		$$CPI_{10} = 118{,}09 \cdot \left(1 + 0{,}0321\right)^{10} \approx 162.$$
	}
\end{ex}

\begin{ex}%[2D4C3-2]%[TT lần 1, Sở GD-ĐT Hưng Yên, 2526]%[Huỳnh Xuân Tín]
	Để chuẩn bị cho lễ hội, một đơn vị thi công dựng một cổng chào dạng vòm Parabol có chiều cao $4$ m và chiều rộng chân cổng là $4$ m. Ở chính giữa cổng, người ta thiết kế một lối đi hình chữ nhật cao $2$ m và rộng $2$ m. Phần diện tích mặt trước của cổng Parabol (sau khi đã trừ lối đi hình chữ nhật) được chia thành $8$ lớp ngang, mỗi lớp cao $0{,}5$ m để ốp tấm nhôm Alu trang trí có màu khác nhau (như bản vẽ thiết kế hình bēn dưới).
	\begin{center}
		\begin{tikzpicture}[line cap=butt,line join=miter,>=stealth]
			\tikzset{declare function={xmin=-1.5;xmax=3.5;ymin=-4.5;ymax=2;
					f(\x)=-0.5*(\x)^2+8;
				},
				smooth,samples=450
			}
			\draw[->] (-5,0)--(5,0) node[shift={(-100:7pt)},font=\normalsize]{$ x $};
			%	\draw[->] (0,ymin)--(0,ymax) node[shift={(190:7pt)},font=\normalsize]{$ y $};
			%	\fill (0,0) node[shift={(225:7pt)},font=\normalsize]{$ O $};
			\fill (-1,0) circle (1pt); \fill (1,-2) circle (1pt); \fill (2,0) circle (1pt); \fill (3,-4) circle (1pt); \fill (0,-4) circle (1pt);
			\draw (0,7)node[above]{Lớp $8$} (0,6)node[above]{Lớp $7$} 
			(0,5)node[above]{Lớp $6$} 
			(0,4)node[above]{Lớp $5$} 
			(-2,0.5)node[left]{Lớp $1$}  	
			(-2,1.5)node[left]{Lớp $2$} 
			(-2,2.5)node[left]{Lớp $3$} 
			(-2,3.5)node[left]{Lớp $4$} 
			(0,-0.3)node[below]{$2$ m} 
			(0,-0.8)node[below]{$4$ m}
			(0,-1.4)node[below]{Chiều rộng} 
			(-5.5,4)node[left]{$4$m} 
			(-4.5,0.5)node[left]{$0{,}5$m} 
			;
			%	\clip (xmin,ymin) rectangle (xmax,ymax);
			\draw  plot[domain=-4:4] (\x, {f(\x)});
			\draw  (-2,0)--(-2,4)--(2,4)--(2,0) (-3.741657,1)--(-2,1) (3.741657,1)--(2,1) (-3.464101615,2)--(-2,2) (3.464101615,2)--(2,2) (-3.16227766,3)--(-2,3) (3.16227766,3)--(2,3) (-2.828427125,4)--(2.828427125,4) (-2.449489743,5)--(2.449489743,5) (-2,6)--(2,6) (-1.41421562,7)--(1.414213562,7);	
			\draw [<->] (-2,-0.3)--(2,-0.3);
			\draw [<->] (-4,-0.8)--(4,-0.8);
			\draw [<->] (-5.7,0)--(-5.7,8);
			\draw [<->] (-4.3,0)--(-4.3,1);
		\end{tikzpicture}	
	\end{center}
	Đơn giá thi công lắp tấm nhôm Alu cho cồng được tính như sau
	\begin{itemize}
		\item  Lớp dưới cùng có giá là $400$ nghìn đồng/m$^2$.
		\item  Càng lên cao, giá thi công mỗi lớp tiếp theo tăng thêm $50$ nghìn đồng/m$^2$ so với lớp ngay dưới nó. 
	\end{itemize}
	Tính tổng chi phí (đơn vị: nghìn đồng) ốp Alu cho phần diện tích còn lại của cổng sau khi đã loại bỏ lối đi (\textit{không làm tròn kết quá các phép toán trung gian, chỉ làm tròn kết quá phép toán cuối cùng đến hàng đơn vị})?
	\shortans[oly]{3814}
	\loigiai{Gắn hệ trục tọa độ $Oxy$ sao cho đỉnh Parabol nằm trên trục $Oy$, trục $Ox$ đi qua chân vòm. Phương trình Parabol là $(P)\colon y=-x^2+4$.
		\begin{center}
			\begin{tikzpicture}[line cap=butt,line join=miter,>=stealth]
				\tikzset{declare function={xmin=-1.5;xmax=3.5;ymin=-4.5;ymax=2;
						f(\x)=-0.5*(\x)^2+8;
					},
					smooth,samples=450
				}
				\draw[->] (-5,0)--(5,0) node[shift={(-100:7pt)},font=\normalsize]{$ x $};
				\draw[->] (0,-1)--(0,9) node[shift={(190:7pt)},font=\normalsize]{$ y $};
				\fill (0,0) node[shift={(225:7pt)},font=\normalsize]{$ O $};
				\fill (-4,0) circle (1pt)node[below]{$-4$};
				\fill (4,0) circle (1pt)node[below]{$4$};
				\fill (-2,0) circle (1pt)node[below]{$-2$};
				\fill (2,0) circle (1pt)node[below]{$2$};			
				\fill (0,8) circle (1pt)node[above right]{$4$};		
				\clip (-6,-1) rectangle (6,9);
				\draw  plot[domain=-4:4] (\x, {f(\x)});
				\draw  (-2,0)--(-2,4)--(2,4)--(2,0) (-3.741657,1)--(-2,1) (3.741657,1)--(2,1) (-3.464101615,2)--(-2,2) (3.464101615,2)--(2,2) (-3.16227766,3)--(-2,3) (3.16227766,3)--(2,3) (-2.828427125,4)--(2.828427125,4) (-2.449489743,5)--(2.449489743,5) (-2,6)--(2,6) (-1.41421562,7)--(1.414213562,7);	
			\end{tikzpicture}	
		\end{center}
		Diện tích chỏm Parabol được tính nhanh bằng công thức $S=\dfrac{2}{3} B h$ (trong dó $B$ là độ rộng đáy, $h$ là chiều cao).\\ 
		Gọi $S_0, S_1, \ldots, S_8$ là diện tích hình phẳng giới hạn bởi $(P)$ và các đường thẳng $y=y_i$. \\
		Chiều cao từ đỉnh Parabol xuống đường $y_i$ là $h=4-y_i$. \\
		Từ phương trình $(P) \Rightarrow x=\sqrt{4-y} \Rightarrow B=2 x=2 \sqrt{4-y_i}$. \\
		Suy ra diện tích chỏm phía trên đường $y=y_i$ là	
		$$S_i=\dfrac{2}{3} \cdot 2 \sqrt{4-y_i} \cdot\left(4-y_i\right)=\dfrac{4}{3}\left(\sqrt{4-y_i}\right)^3.
		$$
		Ta có $S_0(y=0)$; $S_1(y=0,5)$; $S_2(y=1)$; $\ldots$; $S_8(y=4)$. \\
		Diện tích các lớp ngang (chưa trừ phần lõi) là $L_1=S_0-S_1$; $L_2=S_1-S_2$; $\ldots$; $L_8=S_7-S_8$.\\
		Tổng chi phí cho toàn bộ diện tích Parabol (giả sử chưa trừ lõi) là
		\begin{eqnarray*}
			T_1	&=& 400 L_1+450 L_2+\ldots+750 L_8 \\
			&=&400\cdot\left(S_0-S_1\right)+450\cdot\left(S_1-S_2\right)+\ldots+750\cdot\left(S_7-S_8\right)\\
			&=&400\cdot S_0+50\cdot\left(S_1+S_2+\ldots+S_7\right) \\
			&=&400 \cdot \dfrac{4}{3} \cdot\left(\sqrt{4}\right)^3+50 \displaystyle\sum\limits_{y=0{,}5 \text { step } 0{,}5}^{3{,}5} \dfrac{4}{3}\left(\sqrt{4-y}\right)^3.
		\end{eqnarray*}
		Chi phí cho phần lõi (hình chữ nhât $2$ m $\times 2$ m, nằm gọn trong các lớp $L_1$ đến $L_4$, mỗi lớp phần lõi có diện tích $2 \times 0{,}5=1$ m$^2$)
		$$T_2=400 \cdot 1+450 \cdot 1+500 \cdot 1+550 \cdot 1=1\,900.$$
		Vậy tổng chi phí thực tế là
		$$T=T_1-T_2 \approx 3\,814\, \text {(nghìn đồng)}.$$	
	}
\end{ex}


\begin{ex}%[2D1V4-3]%[TT lần 1, Sở GD-ĐT Hưng Yên, 2526]%[Huỳnh Xuân Tín]%[GVPB: Lê Minh]
	Một quần thể vi khuẩn được nuôi cấy trong phòng thí nghiệm. Nồng độ dinh dưỡng $S$ (đơn vị mg/l) thay đổi theo thời gian $t$ giờ $(t \geq 0)$ được mô hình hóa bởi hàm số $S(t)=\dfrac{10 t+5}{t+1}$. Biết rằng tốc độ sinh trưởng $V$ của vi khuẩn phụ thuộc vào nồng độ dinh dưỡng theo hàm số $V(S)=\dfrac{5 S}{S+2}$. Khi thời gian $t$ kéo dài, tốc độ sinh trưởng $V$ tăng dần và ổn định quanh một ngưỡng $K$ nhất định. Hỏi sau bao nhiêu \textbf{phút} thì tốc độ sinh trưởng của vi khuẩn đạt $90\%$ của ngưỡng $K$ đó?
	\shortans[oly]{15}
	\loigiai{Khi thời gian $t$ kéo dài thì ta có 
		\[\lim\limits_{t\to +\infty}S(t)=\lim\limits_{t\to +\infty}\dfrac{10 t+5}{t+1}=10.\]
		Do đó $\lim\limits_{t\to +\infty}V(t)=\lim\limits_{S\to 10}V(S)=\lim\limits_{S\to 10}\dfrac{5 S}{S+2}=\dfrac{5\cdot10}{10+2}=\dfrac{25}{6}.$	 Suy ra $K=\dfrac{25}{6}$.	\\
		Ta có $V(S)=90\%\cdot\dfrac{25}{6}=\dfrac{15}{4}$.\\
		Khi đó $\dfrac{5 S}{S+2}=\dfrac{15}{4}\Leftrightarrow S=6$.\\
		Suy ra $S=6\Leftrightarrow\dfrac{10 t+5}{t+1}=6\Leftrightarrow 10t+5=6t+6\Leftrightarrow t=\dfrac{1}{4}=0{,}25$ h $=15$ phút.
	}
\end{ex}

\begin{ex}%[2H5V1-3]%[Dự án đề TT Toán NH25-26 - Đợt 4 -Lê Hoàng Hạc ]%[Sở Giáo Dục Hưng Yên]%[GVPB: Lê Minh]
	Trong không gian với hệ trục tọa độ $Oxyz$, cho mặt phẳng $(P)$ đi qua điểm $M(2;4;8)$ và cắt các tia $Ox$, $Oy$, $Oz$ lần lượt tại $A(a;0;0)$, $B(0;b;0)$, $C(0;0;c)$ sao cho $OA = 2OB = 4OC$. Tính $T = a+b+c$.
	\shortans{73{,}5}
	\loigiai{
		Vì $A$, $B$, $C$ thuộc các tia $Ox$, $Oy$, $Oz$ nên $a > 0$, $b > 0$, $c > 0$. \\
		Theo giả thiết $a = 2b = 4c \Rightarrow b = \dfrac{a}{2}$, $c = \dfrac{a}{4}$.\\
		Phương trình mặt phẳng $(P)$ theo đoạn chắn là $\dfrac{x}{a} + \dfrac{y}{b} + \dfrac{z}{c} = 1 \Leftrightarrow \dfrac{x}{a} + \dfrac{2y}{a} + \dfrac{4z}{a} = 1$.\\
		Thay tọa độ $M(2;4;8)$ vào phương trình trên ta được $\dfrac{2}{a} + \dfrac{8}{a} + \dfrac{32}{a} = 1 \Rightarrow a = 42$.\\
		Suy ra $b = 21$, $c = 10{,}5$.\\
		Vậy $T = a + b + c = 42 + 21 + 10{,}5 = 73{,}5$.
	}
\end{ex}

\begin{ex}%[0D0C2-7]%[Dự án đề TT Toán NH25-26 - Đợt 4 -Lê Hoàng Hạc ]%[Sở Giáo Dục Hưng Yên]%[GVPB: Lê Minh]
	Cho tập hợp $X = \{1;2;3;4;5;6;7;8\}$. Gọi $S$ là tập hợp tất cả các số tự nhiên có $4$ chữ số được lập từ các chữ số thuộc tập $X$. Chọn ngẫu nhiên một số từ tập hợp $S$. Xác suất để chọn được một số chia hết cho $3$ bằng $\dfrac{a}{b}$ (với $a,b \in \mathbb{N}^*$, $\dfrac{a}{b}$ là phân số tối giản). Tính $T=a+b$.
	\shortans{2731}
	\loigiai{
		Số phần tử của tập hợp $S$ là $n(\Omega) = 8^4 = 4096$.\\
		Chia các phần tử của tập $X$ thành $3$ nhóm dựa trên số dư khi chia cho $3$:
		\begin{itemize}
			\item Nhóm chia hết cho $3$: $X_0 = \{3; 6\}$ gồm $2$ phần tử.
			\item Nhóm chia $3$ dư $1$: $X_1 = \{1; 4; 7\}$ gồm $3$ phần tử.
			\item Nhóm chia $3$ dư $2$: $X_2 = \{2; 5; 8\}$ gồm $3$ phần tử.
		\end{itemize}
		Gọi $A$ là biến cố \lq\lq Số được chọn chia hết cho $3$\rq\rq.\\
		Một số chia hết cho $3$ khi và chỉ khi tổng các chữ số của nó chia hết cho $3$.\\
		Giả sử số được lập lấy $x$ chữ số từ $X_0$, $y$ chữ số từ $X_1$ và $z$ chữ số từ $X_2$.\\
		Ta có điều kiện: $\heva{&x + y + z = 4 \\ &(y + 2z) \,\vdots\, 3}$ (với $x, y, z \in \mathbb{N}$ và $0 \le x, y, z \le 4$).\\
		Ta có các trường hợp thỏa mãn như sau
		\begin{itemize}
			\item \textbf{TH1:} $x = 4, y = 0, z = 0$ ($4$ chữ số thuộc $X_0$).
			Số cách lập là $2^4 = 16$ cách.
			\item \textbf{TH2:} $x = 1, y = 3, z = 0$ ($1$ chữ số thuộc $X_0$, $3$ chữ số thuộc $X_1$).
			Số cách lập là $\mathrm{C}_4^1 \times 2^1 \times 3^3 = 216$ cách.
			\item \textbf{TH3:} $x = 1, y = 0, z = 3$ ($1$ chữ số thuộc $X_0$, $3$ chữ số thuộc $X_2$).
			Số cách lập là $\mathrm{C}_4^1 \times 2^1 \times 3^3 = 216$ cách.
			\item \textbf{TH4:} $x = 2, y = 1, z = 1$ ($2$ chữ số thuộc $X_0$, $1$ thuộc $X_1$, $1$ thuộc $X_2$).
			Số cách lập là $\mathrm{C}_4^2 \times \mathrm{C}_2^1 \times 2^2 \times 3^1 \times 3^1 = 432$ cách.
			\item \textbf{TH5:} $x = 0, y = 2, z = 2$ ($2$ chữ số thuộc $X_1$, $2$ chữ số thuộc $X_2$).
			Số cách lập là $\mathrm{C}_4^2 \times 3^2 \times 3^2 = 486$ cách.
		\end{itemize}
		Từ $5$ trường hợp trên, số các số chia hết cho $3$ trong tập $S$ là $$n(A) = 16 + 216 + 216 + 432 + 486 = 1366.$$
		Xác suất của biến cố $A$ là $\mathrm{P}(A) = \dfrac{n(A)}{n(\Omega)} = \dfrac{1366}{4096} = \dfrac{683}{2048}$.\\
		Vì $\dfrac{683}{2048}$ là phân số tối giản nên $a = 683, b = 2048$.\\
		Vậy $T = a + b = 683 + 2048 = 2731$.
	}
\end{ex}
>>>>>>> 33890a9ce1cfbd3079ed36a93e836cf3ef1fde30
\Closesolutionfile{ans}